\pdfoutput=1
\documentclass[11pt]{amsart}

\usepackage[margin=1.1in]{geometry}
\usepackage{amsmath,amssymb,amsthm,mathtools}
\usepackage{array,booktabs,longtable}
\usepackage{enumitem}
\usepackage{microtype}
\usepackage{hyperref}
\usepackage[nameinlink,capitalize]{cleveref}

\hypersetup{colorlinks=true,linkcolor=blue,citecolor=blue,urlcolor=blue}
\raggedbottom
\emergencystretch=2em
\numberwithin{equation}{section}

\theoremstyle{plain}
\newtheorem{theorem}{Theorem}[section]
\newtheorem{proposition}[theorem]{Proposition}
\newtheorem{lemma}[theorem]{Lemma}
\newtheorem{corollary}[theorem]{Corollary}

\theoremstyle{definition}
\newtheorem{definition}[theorem]{Definition}

\theoremstyle{remark}
\newtheorem{remark}[theorem]{Remark}

\newcommand{\R}{\mathbb R}
\newcommand{\Wloc}{\mathcal W}
\newcommand{\Diss}{\mathcal D}
\newcommand{\Err}{\mathcal E}
\newcommand{\EntropyGrad}{\mathfrak G}
\newcommand{\State}{\mathfrak S}
\newcommand{\SDiff}{\mathrm{SDiff}}
\newcommand{\diver}{\operatorname{div}}
\newcommand{\tr}{\operatorname{tr}}
\newcommand{\supp}{\operatorname{supp}}
\newcommand{\esssup}{\operatorname*{ess\,sup}}
\newcommand{\Ploc}{P_{\mathrm{loc}}}
\newcommand{\kappachi}{\kappa_{\chi,R}}

\title[Conditional localized entropy and compactification]
{A Conditional Localized Entropy Functional and State-Space Compactification \\
for Navier--Stokes Singular Branches}

\author{Guillem Duran Ballester}
\email{guillem@fragile.tech}
\date{}

\subjclass[2020]{35Q30, 35B44, 35B40, 76D05}
\keywords{Navier--Stokes equations, singularity formation, local energy estimates, entropy methods, ancient solutions, state-space compactification}

\begin{document}

\begin{abstract}
Assuming the companion Navier--Stokes manuscripts on local Type~I reduction
\cite{BallesterSetup}, residual Type~I exclusion \cite{BallesterResidual},
and Type~II regularity \cite{BallesterTypeII}, we construct a localized
entropy functional \(\Wloc\) on the singular-branch state space supplied by
those works.  The functional combines local enstrophy, a localized
critical-norm velocity budget in \(L^3\), mean-free local pressure in
\(L^{3/2}\), and a
normalized backward adjoint weight adapted to the modulated drift of the
repaired-gauge representation.  Under the imported companion estimates, the
time derivative of \(\Wloc\) decomposes into an explicit nonnegative
dissipation, a cutoff-normalizer term controlled by the local Caccioppoli
estimate, a harmonic-pressure remainder controlled by the localized
pressure decomposition, a modulation correction controlled by the moving
corrected monotonicity, and a Type~I rescaling defect controlled by the
velocity no-escape estimate.  We then identify, in the constrained first
variation of \(\Wloc\), three Euler--Lagrange residuals -- velocity,
pressure, and adjoint -- whose simultaneous vanishing characterizes a
quantitative \emph{zero-gradient balance}.  This balance is matched
row-by-row with the five terminal model classes of the companion
architecture, and entropy-critical limits are assigned to the first
applicable class in the order (1)--(5).  The localized entropy is therefore
a conditional organizing quantity for the companion program, not a
replacement for its local PDE estimates.
\end{abstract}

\maketitle
\tableofcontents

\section{Introduction}

The local study of possible finite-time singularities for the
three-dimensional incompressible Navier--Stokes equations
\begin{equation}\label{eq:NS}
	\partial_t u - \Delta u + (u\cdot\nabla)u + \nabla p = 0,
	\qquad \diver u = 0,
\end{equation}
on \(\R^3\times(T-\delta,T)\) is governed by a persistent tension between
scale-critical compactness and the nonlocal structure of the pressure.  Since
Leray's construction of global weak solutions \cite{Leray1934}, the basic
local framework has been provided by suitable weak solutions and the partial
regularity results of Scheffer, Caffarelli--Kohn--Nirenberg, Lin, and
subsequent authors
\cite{Scheffer1976,CKN1982,Lin1998,GustafsonKangTsai2007}.  In that
framework, the relevant local quantities are the scale-invariant velocity
and pressure functionals, the local energy inequality, and the compactness
mechanisms surviving parabolic rescaling.  The pressure term is
indispensable but nonlocal, and that is precisely what complicates a clean
compactification of singular branches.

The companion manuscripts \cite{BallesterSetup,BallesterResidual,BallesterTypeII}
respond to this difficulty by a detailed local state-space stratification.
The setup paper \cite{BallesterSetup} reduces a singular branch either to a
bounded centered ancient profile produced by Seregin-type extraction, or to
a local Type~II branch.  The residual paper \cite{BallesterResidual} closes
the remaining Type~I residual strata.  The Type~II paper
\cite{BallesterTypeII} proves the repaired-gauge representation, the
localized pressure decomposition, the compact-cylinder Caccioppoli estimate,
the multibubble and cascade exclusions, the scale-collapse and scale-rigid
routing, the windowwise carrier package, the corrected monotonicity, and
the terminal critical-profile decomposition.  Together these papers yield a
finite decision tree.

\subsection*{Purpose of this paper}

The companion papers organize the singularity problem as a decision tree of
named local strata.  Conditional on the estimates and closure theorems
imported from those papers, this paper shows that the same decision tree can
be read off from a \emph{single monotone scalar quantity} -- a localized
entropy \(\Wloc\).  The point is not to bypass the local PDE estimates but
to package them as the error budget of one identity.  A reader who accepts
the companion estimates obtains, from the present paper, a one-quantity
description of the entire stratification.

The motivation comes from two classical sources.  Following Arnold's formal
interpretation of Euler flow as geodesic motion on the group of
volume-preserving diffeomorphisms \cite{Arnold1966,EbinMarsden1970}, one may
regard the convective term as a transport term built into the ambient
geometry.  Following Hamilton and Perelman \cite{Hamilton1982,Perelman2002},
one expects a useful organizing quantity to couple a backward adjoint weight
to the dissipative part of the equation and to penalize failure of
compactification.  In the present setting, however, the functional must be
strictly local and compatible with the actual PDE estimates available in the
companion papers.  We therefore work with local enstrophy, a localized
critical-norm velocity budget, and mean-free local pressure rather than with
a global geometric invariant, and we measure the adjoint concentration via a
Fisher-information component associated with the local backward adjoint
density on a compact cylinder.

\subsection*{What is rigorously proved here}

The paper is self-contained at the level of the entropy formalism, while its
local PDE inputs are the imported companion packages stated in
\Cref{imp:setup-branch-reduction,imp:typeii-analytic-package,imp:residual-typeI-closure,imp:typeii-terminal-closures}.
Specifically, we prove:

\begin{enumerate}[label=\textup{(\arabic*)},leftmargin=2.2em]
	\item the f-equation associated with the normalized backward adjoint
	weight (\cref{prop:f-equation});
	\item a localized monotonicity identity for \(\Wloc\) with explicit
	error decomposition (\cref{thm:local-monotonicity}), with each error term
	tagged to a specific companion theorem;
	\item the exact constrained first variation
	(\cref{prop:principal-part}) in a fixed independent-variable framework;
	\item a conditional matching of the zero-gradient balance with the five
	terminal model classes (\cref{thm:zero-gradient-recovery}), with the
	matching specified by an explicit topology and table
	(\cref{table:gradient-classes-correspondence});
	\item an ordered-classification theorem assigning every entropy-critical
	limit to the first applicable terminal class
	(\cref{thm:five-classes-ordered}).
\end{enumerate}

The local PDE estimates underlying these statements live in the companion
papers; \cref{sec:companion-input} restates the exact inputs used here and
cites them by theorem number.

\subsection*{Main result}

\begin{theorem}[Conditional entropy reduction of the companion program]
\label{thm:intro-main}
Within the singular-branch framework of
\cite{BallesterSetup,BallesterResidual,BallesterTypeII}, there is a
localized entropy functional
\(\Wloc=\Wloc_{\chi,R}[V,P,\rho](s;\sigma)\) with the following properties.
\begin{enumerate}[label=\textup{(\roman*)},leftmargin=2.2em]
	\item \textbf{Monotonicity (\cref{thm:local-monotonicity}).}
	Along every admissible branch state,
	\[
		\Wloc(s_1;\sigma)-\Wloc(s_2;\sigma)
		\ge \int_{s_1}^{s_2}\Diss_{\chi,R}(\tau;\sigma)\,d\tau
		- \Err_{\chi,R}(s_1,s_2;\sigma),
	\]
	with \(\Diss_{\chi,R}\ge 0\) explicit and
	\(\Err_{\chi,R}=\Err^{\mathrm{cut}}+\Err^{\mathrm{prs}}+\Err^{\mathrm{mod}}+\Err^{\mathrm{T1}}\)
	a sum of four explicit terms, each controlled by a named companion
	theorem.

	\item \textbf{First variation (\cref{prop:principal-part}).}
	The constrained first variation of \(\Wloc\) decomposes into three
	displayed Euler--Lagrange residuals
	\(\EntropyGrad_{\chi,R}=
	\bigl(\mathfrak M^{\mathrm{vel}},\mathfrak M^{\mathrm{prs}},\mathfrak M^{\mathrm{adj}}\bigr)\),
	written out in the fixed independent-variable framework in
	\eqref{eq:vel-principal-full}--\eqref{eq:adj-principal-full}.

	\item \textbf{Compactification (\cref{prop:entropy-compactification}).}
	Every admissible singular branch with finite corrected entropy and error
	budget has an entropy-critical subsequential limit.

	\item \textbf{Conditional classification
	(\cref{thm:zero-gradient-recovery,thm:five-classes-ordered}).}
	Using the imported row-entry dictionary, an admissible zero-gradient branch
	state lies in one of five terminal model classes:
	(1) compact single-core, (2) multiple-core, (3) rotational or
	stationary-hull, (4) scale-cascade, (5) diffuse or tail-defect.
	The correspondence between vanishing components of
	\(\bigl(\Err_{\chi,R},\mathfrak M^{\mathrm{vel}},\mathfrak M^{\mathrm{prs}},\mathfrak M^{\mathrm{adj}}\bigr)\)
	and the five classes is recorded in
	\cref{table:gradient-classes-correspondence}.

	\item \textbf{Ordering (\cref{thm:five-classes-ordered}).}  Every
	entropy-critical limit lies in the \emph{first} applicable class in the
	order (1)--(5).

	\item \textbf{Closure (\cref{thm:eliminate-terminal-classes}).}  Each of
	the five terminal model classes is excluded, under the imported
	hypotheses, by the terminal-class exclusion package
	\Cref{imp:terminal-class-exclusion}.
\end{enumerate}
\end{theorem}

\begin{remark}
\Cref{thm:intro-main} is an organizational theorem: the local PDE content
remains in the companion papers.  What is new here is the entropy
formulation, the explicit error decomposition, and the explicit
correspondence between Euler--Lagrange residuals and terminal model classes.
\end{remark}

\subsection*{Relation to the literature}

The local Navier--Stokes background used throughout consists of Leray's weak
solutions \cite{Leray1934}, the suitable weak solution framework and partial
regularity results \cite{Scheffer1976,CKN1982,Lin1998}, classical
critical-space and local theory
\cite{Ladyzhenskaya1969,Temam1977,ConstantinFoias1988,Kato1984,LemarieRieusset2002},
the endpoint regularity theorem of Escauriaza--Seregin--\v Sver\'ak
\cite{ESS2003}, the self-similar rigidity theorem of
Ne\v cas--R\r u\v zi\v cka--\v Sver\'ak \cite{NRS1996}, the Liouville theorem
of Koch--Nadirashvili--Seregin--\v Sver\'ak \cite{KNSS2009}, and the local
Type~I and ancient-solution inputs of Albritton--Barker and Seregin
\cite{AlbrittonBarker2019,Seregin2012}.  The pressure terms in the present
functional are measured in the same mean-free local norms controlled by the
Calder\'on--Zygmund estimates \cite{CalderonZygmund1952,Stein1970} already
used in \cite{BallesterTypeII}.

\subsection*{Structure of the paper}

\Cref{sec:companion-input} records the companion results imported into the
entropy formulation, with theorem-numbered package statements.  \Cref{sec:functional}
defines the local adjoint weight, derives the f-equation, and defines the
localized entropy functional.  \Cref{sec:monotonicity} carries out the
explicit \(\frac{d}{ds}\Wloc\) calculation, displays the four-piece error
decomposition, and proves the localized monotonicity identity.
\Cref{sec:gradient} computes the constrained first variation and displays
the three Euler--Lagrange residuals in the fixed independent-variable
framework.
\Cref{sec:critical} extracts entropy-critical limits, identifies them with
the five terminal model classes via an explicit correspondence table, and
proves the ordered classification.  \Cref{sec:reduction} applies the imported
terminal-class exclusion package to exclude each terminal class under the
stated hypotheses.

\section{Companion Input and the Admissible State Space}
\label{sec:companion-input}

We import companion results as named theorem packages.  Each package below
restates only the hypotheses and conclusions used in the entropy argument and
cites the stable theorem numbers in the companion manuscripts.  Later sections
refer to these imported packages by the labels introduced here, rather than
repeating long source-paper citations.

\begin{theorem}[Imported setup and branch-reduction package]
\label{thm:companion-package}
\label{imp:setup-branch-reduction}
Let \((u,p)\) be a suitable weak Navier--Stokes solution with a hypothetical
singular point.  The setup paper supplies the scale-invariant local
boundedness criteria, the positive concentration token at a singular point,
the local Type~I/Type~II branch dichotomy, the local Seregin extraction in the
Type~I branch, the pointwise Type~I velocity no-escape estimate, and the
direct Liouville closures for the small, stationary, uniformly tight, and
structure-controlled bounded centered ancient Type~I profiles.  These are
\cite[Theorems~4.1, 4.2, 5.7, 6.5, 7.3, 8.4; Lemma~8.7; Theorems~15.7--15.9,
21.4, and~25.5]{BallesterSetup}.
\end{theorem}

\begin{theorem}[Imported represented Type~II analytic package]
\label{imp:typeii-analytic-package}
Every local Type~II branch produced by \Cref{imp:setup-branch-reduction}
admits the repaired-gauge representation used in this paper, the local
pressure decomposition and pressure-normalization comparison, the stability
topology for subsequences, the renormalized local energy inequality, the
compact-cylinder Caccioppoli estimate, the fixed- and moving-cutoff corrected
monotonicity formulae, and the ordered local state-space decomposition.  These
are
\cite[Theorems~3.26, 3.36--3.38, 5.4, 9.11, 9.19, 10.3, 10.4, 10.17, 10.18,
and~10.25; Propositions~3.35 and~5.1; Lemma~5.3]{BallesterTypeII}.
\end{theorem}

\begin{theorem}[Imported Type~II terminal-closure package]
\label{imp:typeii-terminal-closures}
In the represented Type~II state space of \Cref{imp:typeii-analytic-package},
the compact single-core, multibubble, same-point cascade, separated-point,
rough-core, scale-collapse, scale-rigid, windowwise carrier, cost-divergence,
critical-profile, retained compact-test, diffuse-tail, exterior, and final
local Type~II terminal alternatives are either routed to a previous
alternative or excluded.  This is the combined content of
\cite[Theorems~2.19, 4.19, 4.20, 4.45, 5.4, 5.7, 6.7, 6.20, 6.24, 7.130,
8.6, 8.14, 8.38, 8.41, 8.46, 9.42, 9.82, 9.85, 9.89, 9.91, 9.94, 9.103,
9.113, 9.134, 9.139, 10.5, 10.6, 10.38, and~10.39; Corollaries~7.123,
7.128, 9.43, and~9.142; Proposition~4.43]{BallesterTypeII}.
\end{theorem}

\begin{theorem}[Imported residual Type~I closure package]
\label{imp:residual-typeI-closure}
For a residual Type~I branch satisfying the setup normalization and retained
local \(L^3\)-mass hypotheses, the residual state-space stratification closes
the rotational, stationary-hull, affine/parasitic, critical-tail,
diffuse-defect, generic terminal, and remaining residual classes.  In
particular the residual hypothesis left open by
\Cref{imp:setup-branch-reduction} is verified.  This is
\cite[Theorems~1.14, 4.4, 9.27, 9.38, 9.91, 9.98, 9.113, 10.25, 11.2, 11.3,
and~12.1; Propositions~6.10 and~7.1; Remark~7.4; Corollary~12.2]{BallesterResidual}.
\end{theorem}

\begin{theorem}[Imported terminal-class exclusion package]
\label{imp:terminal-class-exclusion}
Under the admissibility, normalization, branch-topology, retained-mass, and
finite-budget hypotheses used in this paper, the five terminal classes are
excluded by the following companion results:
\begin{center}\small
\setlength{\tabcolsep}{4pt}
\begin{tabular}{@{}>{\raggedright\arraybackslash}p{0.18\textwidth}
				>{\raggedright\arraybackslash}p{0.34\textwidth}
				>{\raggedright\arraybackslash}p{0.40\textwidth}@{}}
\toprule
Terminal class & Companion manifestation & Exclusion input \\
\midrule
\(\mathcal C_{\mathrm{core}}\) &
bounded centered Type~I profiles; compact single-core or retained compact-test
Type~II states &
\cite[Theorems~15.7--15.9, 21.4, and~25.5]{BallesterSetup};
\cite[Theorems~2.19, 9.103, and~10.5]{BallesterTypeII} \\
\(\mathcal C_{\mathrm{multi}}\) &
multibubble, same-point cascade, separated-point, and terminal decoupling
states &
\cite[Theorems~4.19, 4.20, 4.45, 8.6, 8.14, and~10.6]{BallesterTypeII} \\
\(\mathcal C_{\mathrm{rot}}\) &
rotational, stationary-hull, and axisymmetric residual branches &
\cite[Theorem~4.4; Propositions~6.10 and~7.1; Remark~7.4; Theorem~10.25]{BallesterResidual} \\
\(\mathcal C_{\mathrm{cascade}}\) &
scale-collapse, scale-rigid, windowwise carrier, and cost-divergence states &
\cite[Theorems~6.7, 6.20, 6.24, 7.130, and~9.42;
Corollaries~7.123, 7.128, and~9.43; Proposition~4.43]{BallesterTypeII} \\
\(\mathcal C_{\mathrm{tail}}\) &
diffuse, critical-tail, exterior, and terminal inactive defects &
\cite[Theorems~9.27, 9.38, 9.91, 9.98, and~9.113]{BallesterResidual};
\cite[Theorems~8.38, 8.41, 8.46, 9.89, 9.91, and~9.94]{BallesterTypeII} \\
\bottomrule
\end{tabular}
\end{center}
Thus each of the terminal classes
\(\mathcal C_{\mathrm{core}}, \mathcal C_{\mathrm{multi}},
\mathcal C_{\mathrm{rot}}, \mathcal C_{\mathrm{cascade}},
\mathcal C_{\mathrm{tail}}\) has no admissible state satisfying these
hypotheses.
\end{theorem}

\begin{definition}[Admissible branch state]
\label{def:admissible-state}
An \emph{admissible branch state} is either:
\begin{enumerate}[label=\textup{(\roman*)},leftmargin=2.2em]
		\item a bounded centered ancient profile
		\((V,P)\) on \(\R^3\times(-\infty,0)\)
		produced by the Type~I Seregin extraction imported in
		\Cref{imp:setup-branch-reduction}; in this
		case we set the modulation coefficients to the centered self-similar
		values \(a(s)\equiv\tfrac12\), \(b(s)\equiv0\), and set
		\(\nu=1\) in the unified equation below; or
		\item a represented local Type~II branch \((V,P,a,b)\) on a compact
		parabolic cylinder \(Q_R=B_R\times(-R^2,0)\) in the sense of the
		represented Type~II analytic package \Cref{imp:typeii-analytic-package},
		satisfying the renormalized
	suitable equation
	\begin{equation}\label{eq:renorm-NS}
		\partial_\tau V+(V\cdot\nabla)V+\nabla P-\nu\Delta V
		= a(\tau)V+a(\tau)(y\cdot\nabla V)+b(\tau)\cdot\nabla V,
		\qquad \diver V=0,
	\end{equation}
		with \(\nu>0\), \(a,b\in L^\infty(-R^2,0)\), the localized pressure
		decomposition, and the renormalized local energy inequality imported in
		\Cref{imp:typeii-analytic-package}.
\end{enumerate}
We write \(\State\) for the collection of admissible branch states.
\end{definition}

\begin{remark}\label{rem:typeI-as-subcase}
Setting \(a\equiv\tfrac12\), \(b\equiv0\), and \(\nu=1\), the renormalized
equation \eqref{eq:renorm-NS} reduces to the centered self-similar equation
\[
	\partial_\tau V+(V\cdot\nabla)V+\nabla P-\Delta V
	= \tfrac12 V+\tfrac12(y\cdot\nabla V),
\]
which is exactly the centered ancient equation produced by Seregin
extraction.  Thus the Type~I and Type~II cases are simultaneous specializations of
\eqref{eq:renorm-NS}, justifying the unified treatment below.  All
calculations of \cref{sec:monotonicity,sec:gradient} are performed on
\eqref{eq:renorm-NS} and apply to both cases by setting \((a,b,\nu)\) to
their respective values.
\end{remark}

The entropy formalism does not introduce a new state space; it is defined on
the state space already produced by
\Cref{imp:setup-branch-reduction,imp:typeii-analytic-package}.  The role
of the entropy is to organize which branch changes are dissipative, which
are merely routing mechanisms, and which terminal regimes remain after all
routing has been exhausted.

\begin{remark}[Regularity convention]
\label{rem:regularity-convention}
The differential identities in \cref{sec:functional,sec:monotonicity,sec:gradient}
are first stated and proved for smooth represented branch states and smooth
positive adjoint weights on compact subcylinders.  For suitable weak or
limiting branch states, the same formulae are used in the integrated or
distributional sense obtained by applying the smooth identities to the
regularized representatives supplied by the companion compactness and local
energy packages, then passing to the branch topology.  The exact residuals
and the four error functionals in \cref{def:errors} record the terms that
must survive this passage to the limit.
\end{remark}

\section{A Localized Entropy Functional}
\label{sec:functional}

Throughout this section, fix an admissible branch state
\((V,P,a,b)\in\State\) represented on \(Q_R=B_R\times(-R^2,0)\), after
restricting a Type~I ancient profile to \(Q_R\) if necessary, and fix a
terminal parameter \(\sigma\in(-R^2,0)\).  For \(s\in(-R^2,\sigma)\) set
\begin{equation}\label{eq:theta}
	\theta=\theta(s;\sigma):=\sigma-s>0.
\end{equation}
Fix a nonnegative, nontrivial cutoff \(\chi\in C_c^\infty(B_R)\); we keep
\(\chi\) and \(R\) fixed throughout the section.  All errors involving
derivatives of \(\chi\) are systematically tracked.

\subsection{The normalized backward adjoint weight}

\begin{definition}[Normalized backward adjoint weight]
\label{def:adjoint-weight}
A \emph{normalized backward adjoint weight} for \((V,P,a,b)\) on
\((s_0,\sigma)\) is a positive function
\(\rho\in C^2(B_R\times(s_0,\sigma))\), \(s_0>-R^2\), satisfying
\begin{equation}\label{eq:adjoint-rho}
	-\partial_s\rho-\Delta\rho-\diver\bigl((a(s)y+b(s)-V)\rho\bigr)
	+\kappachi(s)\rho=0
	\qquad\text{in }B_R\times(s_0,\sigma),
\end{equation}
together with the scalar normalization coefficient
\begin{equation}\label{eq:kappa-def}
	\kappachi(s)
	:=\frac{\displaystyle\int_{B_R}
	\bigl(\Delta\chi^2-(a(s)y+b(s)-V)\cdot\nabla\chi^2\bigr)\rho\,dy}
	{\displaystyle\int_{B_R}\chi^2\rho\,dy},
\end{equation}
	taken on intervals where the denominator is positive, together with the
	normalization
\begin{equation}\label{eq:rho-normalization}
	\int_{B_R}\chi^2(y)\rho(y,s;\sigma)\,dy=1
	\qquad\text{for every }s\in(s_0,\sigma).
\end{equation}
We write
\begin{equation}\label{eq:rho-f}
	\rho(y,s;\sigma)=(4\pi\theta)^{-3/2}e^{-f(y,s;\sigma)},
\end{equation}
which uniquely determines \(f\in C^2(B_R\times(s_0,\sigma))\) on the support
of \(\rho\).
\end{definition}

\begin{remark}
The drift \(a(s)y+b(s)-V(y,s)\) in \eqref{eq:adjoint-rho} matches the drift
of the renormalized equation \eqref{eq:renorm-NS}.  This is the local
analogue of the conjugate heat equation used in the Hamilton--Perelman
\(\mathcal W\)-functional.  In contrast to that setting, equation
\eqref{eq:adjoint-rho} is solved on a compact spatial domain and the
normalization \eqref{eq:rho-normalization} is imposed via the cutoff
\(\chi\).  The scalar coefficient \(\kappachi\) is not an additional branch
modulation; it is the Lagrange multiplier enforcing the weighted mass
constraint.  Existence for fixed terminal datum can be obtained by solving
the unnormalized adjoint equation and dividing by
\(\int\chi^2\rho\), which produces exactly \eqref{eq:kappa-def} on intervals
where this mass is nonzero.
\end{remark}

\begin{lemma}[Weighted adjoint mass balance]
\label{lem:adjoint-rho-conserves-norm}
For every \(\rho\) satisfying \eqref{eq:adjoint-rho},
\begin{equation}\label{eq:adjoint-rho-conserves-norm}
	\frac{d}{ds}\int_{B_R}\chi^2\rho\,dy
	=-\int_{B_R}\bigl(\Delta\chi^2-(a(s)y+b(s)-V)\cdot\nabla\chi^2\bigr)\rho\,dy
	+\kappachi(s)\int_{B_R}\chi^2\rho\,dy,
\end{equation}
because \(\chi\in C_c^\infty(B_R)\).  In particular, with \(\kappachi\)
defined by \eqref{eq:kappa-def}, the normalization
\eqref{eq:rho-normalization} is preserved once it holds at one time.
\end{lemma}

\begin{proof}
Multiply \eqref{eq:adjoint-rho} by \(\chi^2\), integrate, and integrate by
parts.  Using \(\diver V=0\),
\[
\begin{aligned}
	\frac{d}{ds}\int\chi^2\rho
	&=\int\chi^2\bigl(-\Delta\rho-\diver((a y+b-V)\rho)+\kappachi\rho\bigr)\\
	&=-\int\Delta\chi^2\,\rho+\int\nabla\chi^2\cdot(a y+b-V)\rho
	+\kappachi\int\chi^2\rho+\text{boundary terms}.
\end{aligned}
\]
The boundary terms vanish because the support of \(\chi\) is compactly
contained in \(B_R\).  This is \eqref{eq:adjoint-rho-conserves-norm}.
Substituting \eqref{eq:kappa-def} gives zero derivative.
\end{proof}

We now derive the equation satisfied by \(f\).  This is the essential
piece of preparation for the entropy calculation in
\cref{sec:monotonicity}.

\begin{proposition}[The f-equation]
\label{prop:f-equation}
Let \(\rho=(4\pi\theta)^{-3/2}e^{-f}\) satisfy \eqref{eq:adjoint-rho}.
Then on \(\{\rho>0\}\),
\begin{equation}\label{eq:f-equation}
	\partial_s f+\Delta f-|\nabla f|^2+(a(s)y+b(s)-V)\cdot\nabla f
	=3a(s)+\frac{3}{2\theta}-\kappachi(s).
\end{equation}
Equivalently,
\begin{equation}\label{eq:f-equation-rewritten}
	-\partial_s f
	=\Delta f-|\nabla f|^2+(a(s)y+b(s)-V)\cdot\nabla f
	-3a(s)-\frac{3}{2\theta}+\kappachi(s).
\end{equation}
\end{proposition}

\begin{proof}
By \eqref{eq:rho-f}, \(\rho=(4\pi\theta)^{-3/2}e^{-f}\), hence
\[
\begin{aligned}
	\partial_s\rho
	&=\frac{3}{2}(4\pi\theta)^{-5/2}(4\pi)e^{-f}-(4\pi\theta)^{-3/2}e^{-f}\partial_s f\\
	&=\frac{3}{2\theta}\rho-(\partial_s f)\rho,
\end{aligned}
\]
where we used \(\partial_s\theta=-1\), so
\(\partial_s\bigl((4\pi\theta)^{-3/2}\bigr)
=(4\pi\theta)^{-3/2}\cdot\frac{3}{2\theta}\).  Similarly,
\[
\nabla\rho=-(\nabla f)\rho,\qquad
\Delta\rho=\bigl(|\nabla f|^2-\Delta f\bigr)\rho.
\]
Compute the convective term:
\[
\begin{aligned}
	\diver\bigl((a y+b-V)\rho\bigr)
	&=\diver(a y+b-V)\,\rho+(a y+b-V)\cdot\nabla\rho\\
	&=3a(s)\,\rho-(a y+b-V)\cdot(\nabla f)\rho,
\end{aligned}
\]
using \(\diver(a y)=3a\), \(\diver b=0\), and \(\diver V=0\).
Substituting these into \eqref{eq:adjoint-rho} and dividing by \(\rho>0\),
\[
	-\frac{3}{2\theta}+\partial_s f-\bigl(|\nabla f|^2-\Delta f\bigr)
	-3a(s)+(a y+b-V)\cdot\nabla f+\kappachi(s)=0,
\]
which is equivalent to \eqref{eq:f-equation}.
\end{proof}

\begin{remark}\label{rem:f-equation-shape}
\Cref{eq:f-equation} is a Hamilton--Jacobi-type equation with viscosity
\(\Delta f\), nonlinear gradient term \(-|\nabla f|^2\), and modulated drift
\((a y+b-V)\cdot\nabla f\).  In the centered self-similar Type~I case
\(a\equiv\tfrac12\), \(b\equiv0\), \(\nu=1\), the right-hand side becomes
\(\tfrac32+\tfrac{3}{2\theta}-\kappachi\).  The added scalar
\(\kappachi\) is the cutoff-normalization correction; it disappears in the
unlocalized whole-space heat-kernel model and is included in the cutoff error
in the local theory below.
\end{remark}

\subsection{Pressure normalization}

We write
\begin{equation}\label{eq:pressure-tilde}
	\widetilde P(y,s):=P(y,s)-(P)_{B_R}(s),
	\qquad (P)_{B_R}(s):=\frac{1}{|B_R|}\int_{B_R}P(y,s)\,dy
\end{equation}
for the mean-free local pressure.  By the pressure package in
\Cref{imp:typeii-analytic-package}, this normalization differs from the
localized Calder\'on--Zygmund piece \(\Ploc\) only by a function of time.  We
use \(\widetilde P\) inside the entropy integrand because the
Calder\'on--Zygmund continuity bound and the renormalized local energy
inequality both reference \(\widetilde P\).

\subsection{The localized entropy functional}

\begin{definition}[Localized entropy functional]
\label{def:entropy}
Let \((V,P,a,b)\in\State\), let \(\rho\) be a normalized backward adjoint
weight, and let \(f\) be defined by \eqref{eq:rho-f}.  The corresponding
\emph{localized entropy} at time \(s\) relative to terminal time \(\sigma\)
is
\begin{equation}\label{eq:entropy-def}
\begin{aligned}
	\Wloc_{\chi,R}[V,P,\rho](s;\sigma)
	:= \int_{B_R}\chi^2
	\Bigl[
			\theta |\nabla V|^2
		+ \theta |V|^3
		+ \theta |\widetilde P|^{3/2}
		+ \theta |\nabla f|^2
		+ f - 3
	\Bigr]\rho\,dy.
\end{aligned}
\end{equation}
\end{definition}

\begin{lemma}[Parabolic scaling bookkeeping]
\label{lem:scaling-bookkeeping}
\label{rem:scaling}
Let \(z=\lambda y\), \(t=\lambda^2s\), and
\(\sigma_\lambda=\lambda^{-2}\sigma\).  Set
\[
	V_\lambda(y,s)=\lambda V(z,t),\qquad
	P_\lambda(y,s)=\lambda^2P(z,t),\qquad
	\rho_\lambda(y,s)=\lambda^3\rho(z,t),
\]
\[
	\chi_\lambda(y)=\chi(z),\qquad
	f_\lambda(y,s)=f(z,t),\qquad
	\theta_\lambda(s;\sigma_\lambda)=\lambda^{-2}\theta(t;\sigma).
\]
Then \(\rho_\lambda\,dy=\rho\,dz\),
\(\int\chi_\lambda^2\rho_\lambda\,dy=\int\chi^2\rho\,dz\), and
\[
\begin{aligned}
	\theta_\lambda|\nabla_y f_\lambda|^2
	&=\theta|\nabla_z f|^2,\\
	\theta_\lambda|V_\lambda|^3
	&=\lambda\,\theta|V|^3,\\
	\theta_\lambda|P_\lambda|^{3/2}
	&=\lambda\,\theta|P|^{3/2},\\
	\theta_\lambda|\nabla_y V_\lambda|^2
	&=\lambda^2\theta|\nabla_z V|^2 .
\end{aligned}
\]
Thus the adjoint Fisher part and the cutoff normalization are exactly
covariant, while the velocity and pressure terms have the expected local
coercive weights in the renormalized branch variables.  The entropy
identity below uses these terms as a controlled local budget; no standalone
scale-invariance assertion is used in the proof.
\end{lemma}

\begin{remark}[Heat-kernel normalization]
In the unlocalized whole-space model
\(\rho_*(y,s)=(4\pi\theta)^{-3/2}\exp(-|y|^2/(4\theta))\), one has
\(f=|y|^2/(4\theta)\) and
\[
	\int_{\R^3}(\theta|\nabla f|^2+f)\rho_*\,dy=3.
\]
Therefore the constant \(-3\) normalizes the pure adjoint heat-kernel state
to zero when \(V\equiv0\) and \(P\equiv0\).
\end{remark}

\begin{remark}\label{rem:cutoff-track}
The cutoff \(\chi\) and the radius \(R\) are tracked explicitly: we will
show in \cref{sec:monotonicity} that all derivatives of \(\chi\) appear
inside one named error term \(\Err^{\mathrm{cut}}_{\chi,R}\) and are
controlled by the compact-cylinder Caccioppoli estimate
imported in \Cref{imp:typeii-analytic-package}.
\end{remark}

\section{Monotonicity up to Controlled Local Error}
\label{sec:monotonicity}

We now compute \(\frac{d}{ds}\Wloc_{\chi,R}\) along the renormalized
equation \eqref{eq:renorm-NS}, the f-equation \eqref{eq:f-equation}, and
the local pressure decomposition imported in \Cref{imp:typeii-analytic-package}.
To make the
monotonicity statement auditable, we first record an exact identity with no
implicit commutator terms.  The companion papers are used only afterward,
to dominate the exact residual by the four named error classes.

\subsection{Setup of the calculation}

For convenience write
\begin{equation}\label{eq:integrand}
	\mathcal I(y,s):=\theta |\nabla V|^2+\theta |V|^3
	+\theta |\widetilde P|^{3/2}+\theta |\nabla f|^2+f-3,
\end{equation}
so that \(\Wloc_{\chi,R}=\int\chi^2\mathcal I\rho\,dy\).  We use the
following localized transport identity with
\[
	B(y,s):=a(s)y+b(s)-V(y,s),\qquad
	L:=\partial_s-\Delta+B\cdot\nabla,
\]
throughout the monotonicity calculation.

\begin{lemma}[Localized transport identity]
\label{lem:localized-transport-identity}
Let \(\rho\) be a smooth positive solution of \eqref{eq:adjoint-rho}.  For
any smooth scalar density \(F\), define
\begin{align}
	\mathfrak C_{\chi,R}[F](s)
	&:=\int_{B_R}\bigl[
	-F\Delta\chi^2-2\nabla\chi^2\cdot\nabla F
	+F B\cdot\nabla\chi^2\bigr]\rho\,dy,\label{eq:cutoff-functional}\\
	\mathfrak N_{\chi,R}[F](s)
	&:=\int_{B_R}\chi^2\kappachi(s)F\rho\,dy.\label{eq:normalizer-functional}
\end{align}
Then
\begin{equation}\label{eq:transport-identity}
	\frac{d}{ds}\int_{B_R}\chi^2 F\rho\,dy
	=\int_{B_R}\chi^2 L F\,\rho\,dy
	+\mathfrak C_{\chi,R}[F](s)
	+\mathfrak N_{\chi,R}[F](s).
\end{equation}
\end{lemma}

\begin{proof}
Differentiate under the integral and use \(\partial_s\rho\) from
\eqref{eq:adjoint-rho}:
\[
\begin{aligned}
	\frac{d}{ds}\int\chi^2F\rho
	&=\int\chi^2(\partial_sF)\rho
	+\int\chi^2F\bigl(-\Delta\rho-\diver(B\rho)+\kappachi\rho\bigr)\\
	&=\int\chi^2(\partial_sF)\rho\,dy
	+\int\bigl(-\Delta(\chi^2F)+B\cdot\nabla(\chi^2F)\bigr)\rho\,dy
	+\int\chi^2\kappachi F\rho\,dy,
\end{aligned}
\]
where the boundary integrals vanish by \(\chi\in C_c^\infty(B_R)\);
expanding derivatives of \(\chi^2\) gives
\eqref{eq:transport-identity}.
\end{proof}

Applying \cref{lem:localized-transport-identity} to \(F=\mathcal I\) gives
the exact master identity
\begin{equation}\label{eq:dW-ds-master}
	\frac{d}{ds}\Wloc_{\chi,R}
	=\int_{B_R}\chi^2 L\mathcal I\,\rho\,dy
	+\mathfrak C_{\chi,R}[\mathcal I](s)
	+\mathfrak N_{\chi,R}[\mathcal I](s).
\end{equation}

\subsection{Exact residual and dissipation}

The nonnegative part of the intrinsic calculation is the following
dissipation:
\begin{equation}\label{eq:dissipation}
\boxed{
\begin{aligned}
	\Diss_{\chi,R}(s;\sigma)
	:=
	&\int\chi^2\Bigl[2\nu\theta|\nabla^2 V|^2
	+|\nabla V|^2
	+6\nu\theta|V||\nabla V|^2
	+|V|^3\\
	&\qquad\quad+|\widetilde P|^{3/2}
	+2\theta|\nabla^2 f|^2\Bigr]\rho\,dy.
\end{aligned}}
\end{equation}
The boxed expression is manifestly nonnegative.  No spectral-gap or
Bakry--\'Emery Poincar\'e inequality is used in this version of the entropy
identity.

We now define the exact intrinsic residual by
\begin{equation}\label{eq:intrinsic-residual}
	\mathfrak R_{\chi,R}^{\mathrm{int}}(s;\sigma)
	:=\int_{B_R}\chi^2 L\mathcal I\,\rho\,dy
	+\Diss_{\chi,R}(s;\sigma).
\end{equation}
Equivalently, splitting
\(\mathcal I=\mathcal I_{\nabla V}+\mathcal I_V+\mathcal I_P+\mathcal I_f+\mathcal I_0\)
with
\[
	\mathcal I_{\nabla V}:=\theta|\nabla V|^2,\quad
	\mathcal I_V:=\theta|V|^3,\quad
	\mathcal I_P:=\theta|\widetilde P|^{3/2},\quad
	\mathcal I_f:=\theta|\nabla f|^2,\quad
	\mathcal I_0:=f-3,
\]
the residual is the sum of the five displayed scalar defects
\begin{align}
	\mathfrak R_{\nabla V}
	&:=\int\chi^2 L\mathcal I_{\nabla V}\rho
	+2\nu\theta\int\chi^2|\nabla^2 V|^2\rho
	+\int\chi^2|\nabla V|^2\rho,\label{eq:R-gradV}\\
	\mathfrak R_V
	&:=\int\chi^2 L\mathcal I_V\rho
	+6\nu\theta\int\chi^2|V||\nabla V|^2\rho
	+\int\chi^2|V|^3\rho,\label{eq:R-V}\\
	\mathfrak R_P
	&:=\int\chi^2 L\mathcal I_P\rho
	+\int\chi^2|\widetilde P|^{3/2}\rho,\label{eq:R-P}\\
	\mathfrak R_f
	&:=\int\chi^2 L\mathcal I_f\rho
	+2\theta\int\chi^2|\nabla^2f|^2\rho,\label{eq:R-f}\\
	\mathfrak R_0
	&:=\int\chi^2 L\mathcal I_0\rho.\label{eq:R-zero}
\end{align}
Thus
\[
	\mathfrak R_{\chi,R}^{\mathrm{int}}
	=\mathfrak R_{\nabla V}+\mathfrak R_V+\mathfrak R_P+\mathfrak R_f+\mathfrak R_0.
\]
Every term not explicitly present in \(\Diss_{\chi,R}\), including pressure
couplings, drift/modulation couplings, first-order Fisher terms, and scalar
normalization terms, is therefore visible in one of
\eqref{eq:R-gradV}--\eqref{eq:R-zero}.  Combining
\eqref{eq:dW-ds-master} and \eqref{eq:intrinsic-residual} gives the exact
differential identity
\begin{equation}\label{eq:dW-ds-assembled}
	\frac{d}{ds}\Wloc_{\chi,R}
	=-\Diss_{\chi,R}(s;\sigma)
	+\mathfrak R_{\chi,R}^{\mathrm{int}}(s;\sigma)
	+\mathfrak C_{\chi,R}[\mathcal I](s)
	+\mathfrak N_{\chi,R}[\mathcal I](s).
\end{equation}
This is the identity used below.  All terms outside the displayed
dissipation are accounted for by the exact residual
\(\mathfrak R_{\chi,R}^{\mathrm{int}}\) and the two explicit localization
functionals \(\mathfrak C_{\chi,R}\) and \(\mathfrak N_{\chi,R}\).

\subsection{The four-piece error budget}

We now define the four named error contributions as a domination package for
the exact residual in \eqref{eq:dW-ds-assembled}.  This is the precise place
where the companion papers enter.

\begin{definition}[Named error terms]
\label{def:errors}
On \((s_1,s_2)\subset(-R^2,\sigma)\), a four-piece error decomposition is a
quadruple of nonnegative absolutely continuous interval functionals
\[
	\Err^{\mathrm{cut}}_{\chi,R},\qquad
	\Err^{\mathrm{prs}}_{\chi,R},\qquad
	\Err^{\mathrm{mod}}_{\chi,R},\qquad
	\Err^{\mathrm{T1}}_{\chi,R},
\]
with \(\Err^{\mathrm{T1}}_{\chi,R}\equiv0\) in the Type~II case, such that
there are nonnegative densities
\(\mathfrak e^{\mathrm{cut}},\mathfrak e^{\mathrm{prs}},
\mathfrak e^{\mathrm{mod}},\mathfrak e^{\mathrm{T1}}\in L^1_{\mathrm{loc}}\)
with
\[
	\Err^\bullet_{\chi,R}(s_1,s_2)
	=\int_{s_1}^{s_2}\mathfrak e^\bullet_{\chi,R}(s)\,ds,
\]
and
\begin{equation}\label{eq:error-domination}
\begin{aligned}
	&\int_{s_1}^{s_2}
	\Bigl|
	\mathfrak R_{\chi,R}^{\mathrm{int}}(s;\sigma)
	+\mathfrak C_{\chi,R}[\mathcal I](s)
	+\mathfrak N_{\chi,R}[\mathcal I](s)
	\Bigr|\,ds\\
	&\qquad\le
	\Err^{\mathrm{cut}}_{\chi,R}(s_1,s_2)
	+\Err^{\mathrm{prs}}_{\chi,R}(s_1,s_2)
	+\Err^{\mathrm{mod}}_{\chi,R}(s_1,s_2)
	+\Err^{\mathrm{T1}}_{\chi,R}(s_1,s_2).
\end{aligned}
\end{equation}
The total error is
\[
	\Err_{\chi,R}:=\Err^{\mathrm{cut}}_{\chi,R}
	+\Err^{\mathrm{prs}}_{\chi,R}
	+\Err^{\mathrm{mod}}_{\chi,R}
	+\Err^{\mathrm{T1}}_{\chi,R}.
\]
\end{definition}

\begin{remark}\label{rem:err-interpretation}
\(\Err^{\mathrm{cut}}\) measures derivatives falling on \(\chi\) and the
cutoff normalizer \(\kappachi\), i.e. the terms represented by
\(\mathfrak C_{\chi,R}\) and \(\mathfrak N_{\chi,R}\).
\(\Err^{\mathrm{prs}}\) measures the components of
\(\mathfrak R_{\chi,R}^{\mathrm{int}}\) involving the local pressure
decomposition and its harmonic remainder.
\(\Err^{\mathrm{mod}}\) measures the components of
\(\mathfrak R_{\chi,R}^{\mathrm{int}}\) carrying \(a(s)\), \(b(s)\), and
moving-cutoff modulation.
\(\Err^{\mathrm{T1}}\) measures the Type~I rescaling defect for centered
ancient profiles only.  Each term will now be controlled by a named
companion theorem.
\end{remark}

\subsection{Companion-theorem control of the four error terms}

\begin{theorem}[Imported companion domination of the exact residual]
\label{ass:companion-error-domination}
\label{imp:exact-error-domination}
For every admissible represented branch, every cutoff
\(\chi\in C_c^\infty(B_R)\) with \(\chi\equiv1\) on \(B_{R/2}\), and every
interval \((s_1,s_2)\Subset(-R^2,\sigma)\), the companion estimates provide
a four-piece error decomposition in the sense of \cref{def:errors}.  More
precisely:
\begin{enumerate}[label=\textup{(\roman*)},leftmargin=2.2em]
	\item The cutoff and normalizer terms
	\(\mathfrak C_{\chi,R}[\mathcal I]\) and
	\(\mathfrak N_{\chi,R}[\mathcal I]\) are included in
	\(\Err^{\mathrm{cut}}_{\chi,R}\) and are controlled by the compact
	cylinder Caccioppoli, cutoff, and renormalized local-energy estimates
	\cite[Theorem~5.4 and Proposition~5.1]{BallesterTypeII}.

	\item The pressure components of
	\(\mathfrak R_{\chi,R}^{\mathrm{int}}\), including the harmonic-pressure
	remainder from the local decomposition of \(\widetilde P\), are included
	in \(\Err^{\mathrm{prs}}_{\chi,R}\) and are controlled by the pressure
	decomposition and pressure-normalization comparison
	\cite[Theorem~3.26 and Lemma~5.3]{BallesterTypeII}.

	\item The terms in \(\mathfrak R_{\chi,R}^{\mathrm{int}}\) containing
	the modulation coefficients \(a,b\) and the moving cutoff are included in
	\(\Err^{\mathrm{mod}}_{\chi,R}\) and are controlled by the fixed- and
	moving-cutoff corrected monotonicity formulae
	\cite[Theorems~9.11 and~9.19]{BallesterTypeII}.

	\item In the Type~I rescaling case, the rescaling defect is included in
	\(\Err^{\mathrm{T1}}_{\chi,R}\) and is controlled by the velocity
	no-escape estimate \cite[Theorem~6.5]{BallesterSetup}; in the
	Type~II case this term is identically zero.
\end{enumerate}
Together these estimates imply the domination inequality
\eqref{eq:error-domination}.  This theorem is not an additional PDE
claim of the entropy paper; it records the precise imported content from the
companion manuscripts needed to turn the exact identity
\eqref{eq:dW-ds-assembled} into monotonicity.
\end{theorem}

\begin{lemma}[Assembly of the exact domination inequality]
\label{lem:exact-domination-assembly}
Assume the four imported bounds in \Cref{imp:exact-error-domination}.  Then
the single domination estimate \eqref{eq:error-domination} follows with no
omitted remainder.
\end{lemma}

\begin{proof}
The exact scalar integrand in \eqref{eq:error-domination} is first expanded
using the displayed identity
\(\mathfrak R_{\chi,R}^{\mathrm{int}}
=\mathfrak R_{\nabla V}+\mathfrak R_V+\mathfrak R_P+\mathfrak R_f
+\mathfrak R_0\), together with the explicit cutoff-normalizer terms
\(\mathfrak C_{\chi,R}[\mathcal I]\) and
\(\mathfrak N_{\chi,R}[\mathcal I]\).  The companion estimates quoted in
\Cref{imp:exact-error-domination}\textup{(i)--(iv)} are assumed precisely for
this expanded list of terms: cutoff and normalizer contributions are assigned
to \(\Err^{\mathrm{cut}}\), pressure-decomposition contributions to
\(\Err^{\mathrm{prs}}\), repaired-gauge and moving-window contributions to
\(\Err^{\mathrm{mod}}\), and the Type~I rescaling defect to
\(\Err^{\mathrm{T1}}\).  Applying those four estimates term by term and then
using the triangle inequality in time gives \eqref{eq:error-domination}.  No
spatial boundary term remains because the transport identity is tested with
\(\chi\in C_c^\infty(B_R)\).  This lemma concerns the monotonicity identity
only; the cutoff derivatives appearing in the separate first-variation
identity are tracked separately in \(\mathcal B_{\chi,R}\) and enter later
through the imported zero-gradient bridge.
\end{proof}

\subsection{The localized monotonicity identity}

We can now state the main identity.

\begin{theorem}[Localized monotonicity identity]
\label{thm:local-monotonicity}
Let \((V,P,a,b)\in\State\) be an admissible represented branch on \(Q_R\),
let \(\chi\in C_c^\infty(B_R)\) satisfy \(\chi\equiv1\) on \(B_{R/2}\), and
let \(-R^2<s_1<s_2<\sigma<0\).  If \(\rho\) is a normalized backward
adjoint weight on a time interval containing \([s_1,s_2]\), then
	\begin{equation}\label{eq:monotonicity}
		\Wloc_{\chi,R}(s_1;\sigma)-\Wloc_{\chi,R}(s_2;\sigma)
		\ge
		\int_{s_1}^{s_2}\Diss_{\chi,R}(\tau;\sigma)\,d\tau
		-\Err_{\chi,R}(s_1,s_2;\sigma),
	\end{equation}
where \(\Diss_{\chi,R}\ge 0\) is given by \eqref{eq:dissipation}, and
	\(\Err_{\chi,R}=\Err^{\mathrm{cut}}+\Err^{\mathrm{prs}}+\Err^{\mathrm{mod}}+\Err^{\mathrm{T1}}\)
is the four-piece error of \cref{def:errors}, supplied by
\cref{imp:exact-error-domination}.
\end{theorem}

\begin{proof}
Integrating the exact identity \eqref{eq:dW-ds-assembled} on
\((s_1,s_2)\) gives
\[
	\Wloc_{\chi,R}(s_2;\sigma)-\Wloc_{\chi,R}(s_1;\sigma)
	=-\int_{s_1}^{s_2}\Diss_{\chi,R}(\tau;\sigma)\,d\tau
	+\int_{s_1}^{s_2}
	\bigl[
	\mathfrak R_{\chi,R}^{\mathrm{int}}
	+\mathfrak C_{\chi,R}[\mathcal I]
	+\mathfrak N_{\chi,R}[\mathcal I]
	\bigr]\,d\tau .
\]
By \cref{lem:exact-domination-assembly}, the absolute value of the last
integral is bounded by \(\Err_{\chi,R}(s_1,s_2;\sigma)\).  Rearranging
gives \eqref{eq:monotonicity}.  Nonnegativity of
\(\Diss_{\chi,R}\) is immediate from \eqref{eq:dissipation}.
\end{proof}

\begin{corollary}[Existence of one-sided entropy limits]
\label{cor:entropy-limit}
For every admissible singular branch and every fixed local cutoff
\(\chi\), on every interval where \(\Wloc_{\chi,R}(\cdot;\sigma)\) is finite
and the four error terms are absolutely continuous and additive in their
upper limit, the improper tail error
\[
	\Err_{\chi,R}(s,\sigma;\sigma)
	:=\lim_{r\uparrow\sigma}\Err_{\chi,R}(s,r;\sigma)
\]
exists whenever the total error is finite.  On such intervals, the function
\[
	s\longmapsto
	\Wloc_{\chi,R}(s;\sigma)+\Err_{\chi,R}(s,\sigma;\sigma)
\]
is nonincreasing as \(s\uparrow\sigma\).  Consequently, if the corrected
entropy and terminal error remain finite up to \(\sigma\), then
\(\Wloc_{\chi,R}(\cdot;\sigma)\) has a finite one-sided limit after
accounting for the controlled terminal error.
\end{corollary}

\begin{proof}
Fix \(s_1<s_2<\sigma\).  By additivity,
\[
	\Err_{\chi,R}(s_1,\sigma;\sigma)
	=\Err_{\chi,R}(s_1,s_2;\sigma)
	+\Err_{\chi,R}(s_2,\sigma;\sigma).
\]
Adding this identity to \cref{thm:local-monotonicity} gives
\[
\begin{aligned}
	&\Wloc_{\chi,R}(s_1;\sigma)+\Err_{\chi,R}(s_1,\sigma;\sigma)\\
	&\qquad\ge
	\Wloc_{\chi,R}(s_2;\sigma)+\Err_{\chi,R}(s_2,\sigma;\sigma)
	+\int_{s_1}^{s_2}\Diss_{\chi,R}(\tau;\sigma)\,d\tau .
\end{aligned}
\]
Thus the corrected entropy is finite and monotone decreasing as time
advances toward \(\sigma\).  Monotone convergence yields the asserted
one-sided limit.
\end{proof}

\begin{remark}[Closed-form check on the model heat kernel]
\label{rem:heat-kernel-check}
Let \(V\equiv0\), \(P\equiv0\), \(a\equiv0\), \(b\equiv0\), and
\(\rho_*=(4\pi\theta)^{-3/2}\exp(-|y|^2/(4\theta))\), so
\(f_*=|y|^2/(4\theta)\) and the f-equation \eqref{eq:f-equation} reduces to
\(\partial_s f_*+\Delta f_*-|\nabla f_*|^2=3/(2\theta)\), satisfied
identically.  All velocity and pressure terms vanish, and direct
computation gives
\[
	\theta|\nabla f_*|^2+f_*-3
	=\frac{|y|^2}{4\theta}+\frac{|y|^2}{4\theta}-3
	=\frac{|y|^2}{2\theta}-3.
\]
The Gaussian moments
\[
	\int_{\R^3}\frac{|y|^2}{2\theta}\,
	(4\pi\theta)^{-3/2}e^{-|y|^2/(4\theta)}\,dy=3,
	\qquad
	\int_{\R^3}\rho_*\,dy=1
\]
yield \(\Wloc_{1,\infty}[0,0,\rho_*](s;\sigma)=3-3=0\), confirming that the
constant \(-3\) is the correct dimensional offset.  This is the Perelman
\(\mathcal W\)-style normalization, transported to the local
compact-cylinder setting.
\end{remark}

\section{First Variation and Vanishing Entropy Gradient}
\label{sec:gradient}

To recover the terminal singular geometries from the entropy functional
quantitatively, we use a constrained first variation of \(\Wloc\) as a
diagnostic.  The variational framework is fixed once and for all in this
section: \(V\), \(P\), and \(f\) are varied independently, while the
constraints \(\diver V=0\), \(\int P=0\), and
\(\int\chi^2\rho=1\) are imposed on the test directions.  We do not vary
the Calder\'on--Zygmund pressure map \(P[V]\), nor do we vary the adjoint
equation itself.  Those PDE relations enter later through the branch
equation and the companion row-entry package, not through the first
variation.

\subsection{Constrained perturbations}

Let \((V,P,a,b)\in\State\) be an admissible branch state, fix a cutoff
\(\chi\in C_c^\infty(B_R)\) and a normalized adjoint weight
\(\rho=(4\pi\theta)^{-3/2}e^{-f}\).  We allow simultaneous perturbations
\begin{equation}\label{eq:perturbations}
	V_\varepsilon=V+\varepsilon\varphi,\qquad
	P_\varepsilon=P+\varepsilon\psi,\qquad
	f_\varepsilon=f+\varepsilon\eta,
\end{equation}
with \(\varphi\in C_c^\infty(B_R;\R^3)\) divergence-free,
\(\psi\in C_c^\infty(B_R)\) with \(\int_{B_R}\psi=0\), and
\(\eta\in C_c^\infty(B_R)\) satisfying
\begin{equation}\label{eq:eta-constraint}
	\int_{B_R}\chi^2\eta\rho\,dy=0,
\end{equation}
which preserves the normalization \eqref{eq:rho-normalization} to first
order, since \(\rho_\varepsilon=(4\pi\theta)^{-3/2}e^{-(f+\varepsilon\eta)}\).
We do \emph{not} perturb \((a,b,\nu)\); these are fixed parameters of the
represented branch.

\subsection{The three Euler--Lagrange residuals: full statement}

Recall the integrand
\(\mathcal I=\theta|\nabla V|^2+\theta|V|^3+\theta|\widetilde P|^{3/2}
+\theta|\nabla f|^2+f-3\).

\begin{proposition}[Exact constrained first variation]
\label{prop:principal-part}
For perturbations \eqref{eq:perturbations} as above, the first variation of
\(\Wloc_{\chi,R}\) decomposes as
\begin{equation}\label{eq:first-variation-identity}
	\left.\frac{d}{d\varepsilon}\right|_{\varepsilon=0}
	\Wloc_{\chi,R}[V_\varepsilon,P_\varepsilon,\rho_\varepsilon](s;\sigma)
	= \int_{B_R}\chi^2\rho
	\Bigl(
			\mathfrak M^{\mathrm{vel}}\cdot\varphi
		+ \mathfrak M^{\mathrm{prs}}\,\psi
		+ \mathfrak M^{\mathrm{adj}}\,\eta
	\Bigr)\,dy
	+ \mathcal B_{\chi,R}[\varphi,\psi,\eta],
\end{equation}
where \(\rho_\varepsilon=(4\pi\theta)^{-3/2}e^{-(f+\varepsilon\eta)}\).  The
cutoff contribution is supported where \(\nabla\chi\ne0\):
\begin{equation}\label{eq:cB-explicit}
	\mathcal B_{\chi,R}[\varphi,\psi,\eta]
	=\int_{B_R}
	\bigl[
	-2\theta\,((\nabla V)\nabla\chi^2)\cdot\varphi
	-2\theta\,\eta\,\nabla\chi^2\cdot\nabla f
	\bigr]\rho\,dy,
\end{equation}
and the residual representatives, understood as weighted constrained
functionals and pointwise on the set where their displayed weights are
nonzero, are
\begin{align}
	\mathfrak M^{\mathrm{vel}}(y,s)
	&=-2\theta\rho^{-1}\diver(\rho\,\nabla V)+3\theta|V|V,
	\label{eq:vel-principal-full}\\[4pt]
	\mathfrak M^{\mathrm{prs}}(y,s)
	&=\tfrac32\theta|\widetilde P|^{-1/2}\widetilde P
	-\lambda_{\mathrm{prs}}(s)(\chi^2\rho)^{-1},
	\label{eq:prs-principal-full}\\[4pt]
	\mathfrak M^{\mathrm{adj}}(y,s)
	&=-2\theta\rho^{-1}\diver(\rho\,\nabla f)
	+1-\mathcal I-\Lambda(s).
	\label{eq:adj-principal-full}
\end{align}
The factor \( |\widetilde P|^{-1/2}\widetilde P\) is interpreted as \(0\)
on \(\{\widetilde P=0\}\), the continuous derivative of
\(|\widetilde P|^{3/2}\) at the origin.
Here
\begin{equation}\label{eq:lambda-prs-def}
	\lambda_{\mathrm{prs}}(s)
	:=\frac1{|B_R|}\int_{B_R}\chi^2\rho\,
	\tfrac32\theta|\widetilde P|^{-1/2}\widetilde P\,dy
\end{equation}
is the canonical representative of the mean-free pressure projection, and
\(\Lambda(s)\) is the Lagrange multiplier
\begin{equation}\label{eq:Lambda-def}
	\Lambda(s):=\int_{B_R}\chi^2\rho\,
	\bigl[-2\theta\rho^{-1}\diver(\rho\nabla f)+1-\mathcal I\bigr]\,dy
\end{equation}
for the normalized-adjoint constraint \eqref{eq:eta-constraint}.  With these
choices, \(\chi^2\rho\,\mathfrak M^{\mathrm{prs}}\) has zero Lebesgue mean
and \(\mathfrak M^{\mathrm{adj}}\) has zero \(\chi^2\rho\)-mean.
\end{proposition}

\begin{proof}
Differentiate \eqref{eq:entropy-def} in \(\varepsilon\) and evaluate at
\(\varepsilon=0\).  Since this is the independent-variable framework,
\(P\) does not vary through \(V\), and \(\rho\) varies only through
\(f_\varepsilon\).

The \(V\)-terms give
\[
	\int\chi^2\rho\,
	\bigl(2\theta\nabla V:\nabla\varphi+3\theta|V|V\cdot\varphi\bigr)\,dy.
\]
Integrating only the interior weight \(\rho\) by parts gives
\[
	2\theta\int\chi^2\rho\,\nabla V:\nabla\varphi
	=\int\chi^2\rho\,
	\bigl[-2\theta\rho^{-1}\diver(\rho\nabla V)\bigr]\cdot\varphi\,dy
	-2\theta\int\rho\,((\nabla V)\nabla\chi^2)\cdot\varphi\,dy.
\]
Thus the first term in \eqref{eq:cB-explicit} is exactly the cutoff
correction not included in the interior residual
\eqref{eq:vel-principal-full}.  Since \(\varphi\) is divergence-free, the
interior representative is understood modulo gradients.

The pressure term gives
\[
	\int\chi^2\rho\,
	\tfrac32\theta|\widetilde P|^{-1/2}\widetilde P\,\psi\,dy.
\]
Here \(\delta\widetilde P=\psi\) because the direction \(\psi\) has zero
Lebesgue mean.  The same constraint means that subtracting
\(\lambda_{\mathrm{prs}}(\chi^2\rho)^{-1}\) does not change the functional.
The value in \eqref{eq:lambda-prs-def} is chosen only to fix the representative
by imposing zero Lebesgue mean on \(\chi^2\rho\,\mathfrak M^{\mathrm{prs}}\).
This is \eqref{eq:prs-principal-full}.

For the adjoint variable,
\(\partial_\varepsilon\rho_\varepsilon|_{\varepsilon=0}=-\eta\rho\).  Hence
\[
	\delta_f\Wloc
	=\int\chi^2\rho\,
	\bigl(2\theta\nabla f\cdot\nabla\eta+(1-\mathcal I)\eta\bigr)\,dy.
\]
Again integrating only the interior weight \(\rho\) by parts gives
\[
	2\theta\int\chi^2\rho\,\nabla f\cdot\nabla\eta
	=\int\chi^2\rho\,
	\bigl[-2\theta\rho^{-1}\diver(\rho\nabla f)\bigr]\eta\,dy
	-2\theta\int\rho\,\eta\,\nabla\chi^2\cdot\nabla f\,dy.
\]
This gives \eqref{eq:adj-principal-full} together with the second cutoff
term in \eqref{eq:cB-explicit}.  Since
\(\int\chi^2\rho\,dy=1\), the choice \eqref{eq:Lambda-def} subtracts the
weighted mean of the interior adjoint coefficient; subtracting this constant
does not change the functional on test functions satisfying
\(\int\chi^2\eta\rho=0\).
\end{proof}

\begin{remark}\label{rem:no-hidden-terms}
The residuals in \eqref{eq:vel-principal-full}--\eqref{eq:adj-principal-full}
are interior variational residuals of \(\Wloc\), not the Navier--Stokes
equation itself.  All derivatives of the cutoff have been placed in the
single explicit functional \(\mathcal B_{\chi,R}\), so there is no double
counting between \(\mathfrak M^{\mathrm{vel}},\mathfrak M^{\mathrm{adj}}\)
and \(\mathcal B_{\chi,R}\).  The pressure equation, the modulated branch
equation, and the adjoint equation are imposed separately through \(\State\)
and \cref{imp:exact-error-domination}.  This avoids mixing two different
variational problems: independent \((V,P,f)\) variation and reduced
variation through \(P[V]\) or \(\rho[V]\).
\end{remark}

\subsection{Zero-gradient branch states}

\begin{definition}[Zero-gradient branch state]
\label{def:zero-gradient-state}
An admissible branch state \(S=(V,P,a,b)\in\State\) is a
\emph{zero-gradient branch state} if there exist a cutoff
\(\chi\in C_c^\infty(B_R)\), a normalized adjoint weight
\(\rho=(4\pi\theta)^{-3/2}e^{-f}\), a row index
\(k\in\{1,\dots,5\}\), and admissible row data \(\mathcal D_k\), such that
the constrained first variation vanishes after the row normalization:
\begin{equation}\label{eq:zero-gradient-condition}
	\mathcal Z_k(S,\rho,\chi;\mathcal D_k)=0
	\qquad\text{in }\mathcal T_{\mathrm{br}}.
\end{equation}
Here
\[
	\langle A,\phi\rangle_{\chi^2\rho}:=\int_{B_R}\chi^2\rho\,A\phi\,dy
\]
for scalar residuals, with the analogous dot product for vector residuals,
and \(\mathcal Z_k=0\) means the following simultaneous conditions.
\begin{enumerate}[label=\textup{(\roman*)},leftmargin=2.2em]
	\item The pressure residual vanishes in the dual of mean-free pressure
	variations:
	\[
		\langle \mathfrak M^{\mathrm{prs}},\psi\rangle_{\chi^2\rho}=0
		\qquad
		\text{for every }\psi\in C_c^\infty(B_R),\quad
		\int_{B_R}\psi\,dy=0.
	\]

	\item The adjoint residual vanishes in the dual of normalization-preserving
	adjoint variations, except in row \(k=5\), where the only allowed
	nonzero part is the tail measure specified by \(\mathcal D_5\):
	\[
		\langle \mathfrak M^{\mathrm{adj}},\eta\rangle_{\chi^2\rho}=0
		\qquad
		\text{for every }\eta\in C_c^\infty(B_R),\quad
		\int_{B_R}\chi^2\eta\rho\,dy=0,
	\]
	after discarding the row-\(5\) tail defect outside every compact subset
	on which the companion topology has retained mass.

	\item The velocity residual vanishes in the divergence-free dual after the
	row-specific normalization.  For rows \(k=1,2,4\), this says that
	\(\mathfrak M^{\mathrm{vel}}\) is a gradient on the retained compact
	core, on each separated core, or on each selected scale window,
	respectively.  For row \(k=3\), there is a constant
	\(\Omega\in\mathfrak{so}(3)\) in \(\mathcal D_3\) such that
	\[
		\mathfrak M^{\mathrm{vel}}
		-\theta\bigl(\Omega V-(\Omega y)\cdot\nabla V\bigr)
	\]
	is a gradient in the divergence-free dual.  For row \(k=5\), the velocity
	residual is a gradient modulo the tail defect measure in \(\mathcal D_5\).

	\item The cutoff first-variation functional satisfies
	\begin{equation}\label{eq:cB-vanish}
		\mathcal B_{\chi,R}\to0
	\end{equation}
	in \(\mathcal T_{\mathrm{br}}\), after the same row normalization.
\end{enumerate}
For \(k=1\) and no defect data, \eqref{eq:zero-gradient-condition} reduces
to the strict vanishing of
\(\mathfrak M^{\mathrm{vel}},\mathfrak M^{\mathrm{prs}},
\mathfrak M^{\mathrm{adj}}\) and \(\mathcal B_{\chi,R}\) as constrained
functionals.
\end{definition}

\begin{remark}\label{rem:component-vanishing}
Because the three residuals come with the constraints
\(\diver\varphi=0\), \(\int\psi=0\), \(\int\chi^2\eta\rho=0\) on the
admissible perturbations, all equalities above are equalities in constrained
duals.  Thus gradients are invisible to the velocity variation, constants
are invisible to the mean-free pressure variation after the projection
\eqref{eq:lambda-prs-def}, and constants in the adjoint Euler--Lagrange
equation are fixed by the multiplier \(\Lambda(s)\).  Rows (3)--(5) do not
change the Euler--Lagrange residuals; they specify the additional symmetry,
window normalization, or tail defect that the companion topology permits
before testing the same constrained duals.
\end{remark}

\subsection{Quantitative correspondence with the five terminal classes}

\begin{definition}[Limiting branch topology]
\label{def:branch-topology}
The limiting topology used in \cref{table:gradient-classes-correspondence}
is the companion branch topology \(\mathcal T_{\mathrm{br}}\).  It consists
of local strong convergence of \(V\) in \(L^3_{\mathrm{loc}}\), local strong
convergence of \(\widetilde P\) in \(L^{3/2}_{\mathrm{loc}}\) modulo
mean-free normalization, weak convergence of the adjoint measures
\(\rho\,dy\) on compact subsets, convergence of the modulation parameters
\((a,b)\) in the topology supplied by the repaired-gauge representation,
and convergence of the four interval error functionals in
\cref{def:errors}.  The Type~I branch uses the Seregin extraction topology
imported in \Cref{imp:setup-branch-reduction}; the Type~II branch uses the
stability and state-decomposition topology imported in
\Cref{imp:typeii-analytic-package}.
\end{definition}

We now state the imported correspondence used by
\cref{thm:zero-gradient-recovery}.  The table records the residual pattern
for each terminal class; the row-entry implication itself is imported from
the companion state-decomposition theorems and stated explicitly below.  In
the scale-cascade row we distinguish the \emph{raw modulation carrier} from
the normalized error \(\Err^{\mathrm{mod}}\) appearing in
\cref{def:entropy-critical}.

\begin{table}[htbp]
\caption{Quantitative correspondence between zero-gradient branch states
and terminal model classes.  Each row records the residuals required to
vanish, the row-specific modulation or defect datum, and the imported theorem
package that closes the corresponding terminal class.  The notation
\(\Err^{*}\to0\) means the named normalized error term vanishes in the
limiting branch topology.}
\label{table:gradient-classes-correspondence}
\centering\small
\begin{tabular}{@{}>{\raggedright\arraybackslash}p{0.18\textwidth}
								>{\raggedright\arraybackslash}p{0.30\textwidth}
								>{\raggedright\arraybackslash}p{0.20\textwidth}
								>{\raggedright\arraybackslash}p{0.24\textwidth}@{}}
\toprule
Class & Vanishing residuals \& errors & Active modulation & Closing theorem \\
\midrule
(1) Compact single-core
& \(\mathfrak M^{\mathrm{vel}}=\theta\nabla\Pi\),
\(\mathfrak M^{\mathrm{prs}}=0\) in the mean-free dual,
\(\mathfrak M^{\mathrm{adj}}=0\) in the normalized-adjoint dual,
\(\Err^{\mathrm{cut}},\Err^{\mathrm{prs}},\Err^{\mathrm{T1}}\to0\),
\(\Err^{\mathrm{mod}}\to0\)
& \(a=\frac12\), \(b=0\) (Type~I) or \(a,b\to\) const (Type~II)
& \Cref{imp:setup-branch-reduction,imp:typeii-terminal-closures} \\[3pt]
\midrule
(2) Multiple-core
& \(\mathfrak M^{\mathrm{vel}}\) is a sum of localized gradients
\(\sum_j\theta\nabla\Pi_j\) supported near separated cores;
\(\mathfrak M^{\mathrm{prs}}=0\) in the mean-free dual on each core;
\(\mathfrak M^{\mathrm{adj}}=0\) in the normalized-adjoint dual;
\(\Err^{*}\to0\) on each core
& \(a,b\) bounded; per-core modulation
& \Cref{imp:typeii-terminal-closures} \\[3pt]
\midrule
(3) Rotational / stationary-hull
& \(\partial_s V=0\) modulo rotation; in the residual-dual formulation,
\(\mathfrak M^{\mathrm{vel}}
-\theta(\Omega V-(\Omega y)\cdot\nabla V)=\theta\nabla\Pi\) with
\(\Omega\) constant; \(\mathfrak M^{\mathrm{prs}}=0\) in the mean-free dual;
\(\mathfrak M^{\mathrm{adj}}=0\) in the normalized-adjoint dual;
\(\Err^{*}\to0\)
& \(b\) remains the translation modulation; \(\Omega\in\mathfrak{so}(3)\)
is a separate residual symmetry parameter
& \Cref{imp:residual-typeI-closure} \\[3pt]
\midrule
(4) Scale-cascade
& \(\mathfrak M^{\mathrm{vel}}=\theta\nabla\Pi\) holds on each scale
window; \(\mathfrak M^{\mathrm{prs}}=0\) in the mean-free dual;
\(\mathfrak M^{\mathrm{adj}}=0\) in the normalized-adjoint dual;
\(\Err^{\mathrm{cut}},\Err^{\mathrm{prs}},\Err^{\mathrm{T1}}\to0\);
normalized \(\Err^{\mathrm{mod}}\to0\) on selected windows
& raw modulation carrier \(J_1\), \(J_6\), or \(J_7\) before row-specific normalization
& \Cref{imp:typeii-terminal-closures} \\[3pt]
\midrule
(5) Diffuse / tail-defect
& \(\mathfrak M^{\mathrm{vel}}=\theta\nabla\Pi\) modulo a tail defect
measure; \(\mathfrak M^{\mathrm{prs}}=0\) in the mean-free dual;
\(\mathfrak M^{\mathrm{adj}}=0\) modulo a tail concentration in
\(\rho\); \(\Err^{*}\to0\) inside compact, defect supported in tail
& \(a,b\) bounded
& \Cref{imp:residual-typeI-closure,imp:typeii-terminal-closures} \\
\bottomrule
\end{tabular}
\end{table}

\begin{theorem}[Imported entropy-critical zero-gradient bridge]
\label{imp:entropy-critical-zero-gradient}
In the limiting branch topology supplied by
\Cref{imp:setup-branch-reduction,imp:typeii-analytic-package}, every
entropy-critical limit obtained from \Cref{thm:local-monotonicity} is a
zero-gradient branch state in the sense of \cref{def:zero-gradient-state}.
This is the zero-gradient bridge supplied by the ordered local state-space
decomposition and compactness topology
\cite[Theorems~10.17 and~10.18; Proposition~3.35]{BallesterTypeII}, together
with the residual closure bridge
\cite[Theorems~1.14 and~12.1; Corollary~12.2]{BallesterResidual}.
\end{theorem}

\begin{theorem}[Imported zero-gradient row-entry dictionary]
\label{prop:conditional-row-entry-dictionary}
In the limiting branch topology supplied by
\Cref{imp:setup-branch-reduction,imp:typeii-analytic-package}, every
admissible zero-gradient branch state satisfying the finite-budget hypotheses
of the companion state space enters one of the five rows of
\cref{table:gradient-classes-correspondence}; after applying the order
\((1)\)--\((5)\), the row is unique.  The exact companion theorem numbers for
the five rows are displayed in \Cref{imp:terminal-class-exclusion}.  More
precisely:
\begin{enumerate}[label=\textup{(\roman*)},leftmargin=2.2em]
	\item A zero-gradient state with one retained compact core and vanishing
	cutoff, pressure, modulation, and Type~I rescaling errors satisfies the
	compact single-core hypotheses closed by
	\Cref{imp:setup-branch-reduction,imp:typeii-terminal-closures}.

	\item A zero-gradient state whose velocity residual splits into finitely
	many separated localized gradients satisfies the multibubble or
	separated-profile hypotheses closed by
	\Cref{imp:typeii-terminal-closures}.

	\item A zero-gradient state with a residual skew parameter
	\(\Omega\in\mathfrak{so}(3)\), in the sense that the velocity residual is
	a gradient modulo \(\Omega V-(\Omega y)\cdot\nabla V\), satisfies the
	rotational or stationary-hull hypotheses closed by
	\Cref{imp:residual-typeI-closure}.

	\item A zero-gradient state whose compact-window residuals vanish but
	whose raw modulation carrier is carried by the canonical windowwise schedule
	satisfies the scale-cascade hypotheses closed by
	\Cref{imp:typeii-terminal-closures}.

	\item A zero-gradient state whose compact residuals vanish only modulo a
	tail defect satisfies the diffuse or tail-defect hypotheses of
	\Cref{imp:residual-typeI-closure,imp:typeii-terminal-closures}.
\end{enumerate}
\end{theorem}

\begin{remark}
Each item is a direct restatement of the row hypotheses in
\cref{table:gradient-classes-correspondence} together with the imported
theorem package that supplies the corresponding state-space entry.  The
dictionary records the exact external input needed by the entropy argument
and does not replace the companion proofs.
\end{remark}

\begin{theorem}[Imported zero-gradient row recovery]
\label{thm:zero-gradient-recovery}
Let \(S_\infty\in\State\) be an entropy-critical limit.  Then \(S_\infty\) is
a zero-gradient branch state in the sense of \cref{def:zero-gradient-state}
and belongs to a unique ordered row of
\cref{table:gradient-classes-correspondence}.  Conversely, a state satisfying
one of the rows of \cref{table:gradient-classes-correspondence} satisfies the
zero-gradient condition in \(\mathcal T_{\mathrm{br}}\), with the row-specific
normalization in the scale-cascade case and modulo the displayed tail defect
in the diffuse/tail row.
\end{theorem}

\begin{proof}
\Cref{imp:entropy-critical-zero-gradient} gives the zero-gradient conclusion
for entropy-critical limits.  The unique ordered row assignment is exactly
\Cref{prop:conditional-row-entry-dictionary}.  Conversely, fix a row of
\cref{table:gradient-classes-correspondence} and the corresponding limiting
state \(S\).  Each row specifies which residuals vanish.  For every row,
the velocity residual is required to be a gradient (possibly summed over
cores in row (2), with a rotational correction in row (3), with row-specific
scale normalization in row (4), or modulo a tail defect in row (5)); by
\cref{rem:component-vanishing}, this is precisely the divergence-free
constraint saying \(\mathfrak M^{\mathrm{vel}}=0\) modulo the allowed row
defect.  Similarly \(\mathfrak M^{\mathrm{prs}}=0\) in the mean-free dual is
the pressure constraint, and \(\mathfrak M^{\mathrm{adj}}=0\) in the
normalized-adjoint dual, modulo the displayed tail concentration in row (5),
is the normalization-preserving constraint.  Hence \(S\) satisfies
\eqref{eq:zero-gradient-condition} in \(\mathcal T_{\mathrm{br}}\), with the
row-specific convention already stated in the theorem.
The vanishing of \(\mathcal B_{\chi,R}\) in the limit is part of the
imported zero-gradient bridge and row-entry dictionary: the same cutoff
compactness estimates that place
\(\mathfrak C_{\chi,R}[\mathcal I]\) and
\(\mathfrak N_{\chi,R}[\mathcal I]\) in \(\Err^{\mathrm{cut}}\) also give
the dual convergence of the first-variation cutoff functional
\eqref{eq:cB-explicit} in \(\mathcal T_{\mathrm{br}}\).
\end{proof}

\begin{remark}
\Cref{thm:zero-gradient-recovery} relies on the imported zero-gradient bridge
\cref{imp:entropy-critical-zero-gradient} and the imported row-entry package
\cref{prop:conditional-row-entry-dictionary}.  The entropy argument supplies
the exact residuals and the four-piece error pattern; the companion
state-decomposition theorems supply the conversion from that pattern to one of
the five terminal model classes.  This is the intended conditional division of
labor.
\end{remark}

\section{Entropy-Critical Limits and the Ordered Classification}
\label{sec:critical}

\subsection{Entropy-critical limits}

\begin{definition}[Entropy-critical limit]
\label{def:entropy-critical}
An admissible branch state \(S_\infty\in\State\) is \emph{entropy-critical}
if there exist admissible represented branch segments \(S^{(n)}(\tau)\) on a
common normalized cylinder, times \(s_n\uparrow\sigma\), and a fixed cutoff
\(\chi\in C_c^\infty(B_R)\) such that:
\begin{enumerate}[label=\textup{(\roman*)},leftmargin=2.2em]
		\item \(S^{(n)}(s_n)\to S_\infty\) in the companion compactness topology
		imported in
		\Cref{imp:setup-branch-reduction,imp:typeii-analytic-package};
	\item \(\Wloc_{\chi,R}[S^{(n)}](s_n;\sigma)\) converges to a finite limit;
	\item there are window lengths \(\delta_n\downarrow0\), with
	\(s_n-\delta_n>-R^2\), such that the dissipation and error contributions
	of \cref{thm:local-monotonicity} satisfy
	\begin{equation}\label{eq:critical-vanishing}
		\int_{s_n-\delta_n}^{s_n}
		\Diss_{\chi,R}^{S^{(n)}}(\tau;\sigma)\,d\tau\to0,
		\qquad
		\Err^{\bullet,S^{(n)}}_{\chi,R}
		\bigl(s_n-\delta_n,s_n;\sigma\bigr)\to0
	\end{equation}
	for each \(\bullet\in\{\mathrm{cut},\mathrm{prs},\mathrm{mod},\mathrm{T1}\}\),
	with the convention that \(\Err^{\mathrm{T1}}\equiv0\) in the Type~II case.
	The superscript \(S^{(n)}\) indicates that the dissipation and error
	functionals are evaluated along the branch segment \(S^{(n)}(\tau)\).
\end{enumerate}
\end{definition}

\begin{remark}
\Cref{def:entropy-critical} requires the four error components to vanish
individually, not only in total, and only on shrinking terminal windows.  This
is the window strength obtained from a finite integrated error budget.  In
particular, a sequence in the scale-cascade row (4) of the table fails the
definition unless \(\Err^{\mathrm{mod}}\to0\) on the selected windowwise
subsequence supplied by \Cref{imp:typeii-terminal-closures}.
\end{remark}

\begin{lemma}[Shrinking-window selection]
\label{lem:shrinking-window-selection}
Let \(F_0,F_1,\dots,F_m\ge0\) be integrable on \((s_*,\sigma)\).  Then
there are \(s_n\uparrow\sigma\) and \(\delta_n\downarrow0\), with
\(s_n-\delta_n>s_*\), such that
\[
	\int_{s_n-\delta_n}^{s_n}F_j(\tau)\,d\tau\to0
	\qquad\text{for every }j=0,\dots,m.
\]
\end{lemma}

\begin{proof}
Choose any \(\delta_n\downarrow0\) and \(s_n\uparrow\sigma\) with
\((s_n-\delta_n,s_n)\subset(s_*,\sigma)\).  Absolute continuity of the
Lebesgue integral gives the displayed convergence for each fixed \(F_j\).
Since the family is finite, the same sequence works for all \(j=0,\dots,m\).
\end{proof}

\begin{lemma}[Uniform shrinking-window selection]
\label{lem:uniform-shrinking-window-selection}
Let \(s_n\uparrow\sigma\).  For each \(n\), let
\(F_{n,0},\dots,F_{n,m}\ge0\) be nonnegative densities on
\((s_*,\sigma)\).  Assume the finite family is uniformly integrable in
terminal neighborhoods:
\[
	\lim_{\delta\downarrow0}
	\sup_n\sum_{j=0}^m\int_{\sigma-\delta}^{\sigma}F_{n,j}(\tau)\,d\tau=0.
\]
Then, after passing to a subsequence, there are \(\delta_n\downarrow0\) with
\(s_n-\delta_n>s_*\) such that
\[
	\int_{s_n-\delta_n}^{s_n}F_{n,j}(\tau)\,d\tau\to0
	\qquad\text{for every }j=0,\dots,m.
\]
\end{lemma}

\begin{proof}
Choose \(\varepsilon_\ell\downarrow0\) so that
\[
	\sup_n\sum_{j=0}^m
	\int_{\sigma-\varepsilon_\ell}^{\sigma}F_{n,j}(\tau)\,d\tau
	\le \ell^{-1}.
\]
Passing to a subsequence, arrange \(\sigma-s_n<\varepsilon_n/2\).  Choose
\[
	0<\delta_n<\min\{\varepsilon_n/2,\;1/n,\;s_n-s_*\}.
\]
Then \((s_n-\delta_n,s_n)\subset(\sigma-\varepsilon_n,\sigma)\), so each
displayed integral is bounded by the corresponding terminal-neighborhood
sum, which tends to zero.
\end{proof}

\begin{proposition}[Entropy compactification]
\label{prop:entropy-compactification}
Every hypothetical singular branch in the companion state space satisfying
the finite corrected entropy budget
\[
	\sup_{s\in(s_*,\sigma)}\bigl|\Wloc_{\chi,R}(s;\sigma)\bigr|<\infty,
	\qquad
	\lim_{s\uparrow\sigma}\Err_{\chi,R}(s_*,s;\sigma)<\infty
\]
for some \(s_*<\sigma\), with the four error components supplied by
\Cref{imp:exact-error-domination}, has an entropy-critical subsequential
limit.
\end{proposition}

\begin{proof}
Let
\[
	E_*(\sigma):=\lim_{s\uparrow\sigma}\Err_{\chi,R}(s_*,s;\sigma)<\infty.
\]
By additivity, the tail error is
\[
	\Err_{\chi,R}(s,\sigma;\sigma)
	=E_*(\sigma)-\Err_{\chi,R}(s_*,s;\sigma),
\]
and tends to \(0\) as \(s\uparrow\sigma\).  Hence
\cref{cor:entropy-limit} gives a finite one-sided limit for
\(\Wloc_{\chi,R}(\cdot;\sigma)\).
The compactness components of
\Cref{imp:setup-branch-reduction,imp:typeii-analytic-package} provide
subsequential limits in the corresponding state-space topologies.  Telescoping
\begin{equation}\label{eq:telescoping}
	\sum_n\bigl(\Wloc(s_n;\sigma)-\Wloc(s_{n+1};\sigma)\bigr)
	+\sum_n\Err_{\chi,R}(s_n,s_{n+1};\sigma)
	\ge \sum_n\int_{s_n}^{s_{n+1}}\Diss\,ds
\end{equation}
along a sequence \(s_n\uparrow\sigma\), together with the finiteness of
the entropy limit and the finite integrated error budget, gives finite
integrated dissipation and finite integrated error near the terminal time for
this fixed singular branch.  Applying
\cref{lem:shrinking-window-selection} to the dissipation density and the
four error densities gives shrinking terminal windows on which each
component vanishes.  Passing to the subsequence selected by the companion
compactness theorem preserves this vanishing.  If the companion extraction is
formulated instead as a sequence of already rescaled representatives, the
same conclusion follows from \cref{lem:uniform-shrinking-window-selection}
provided the companion compactness package supplies the stated uniform
terminal integrability on the common cylinder.  Without that uniform version
one applies the fixed-branch selection before rescaling, which is the case
used above.  The resulting limit is entropy-critical.
\end{proof}

\subsection{Ordered classification of entropy-critical limits}

The point of the entropy formalism is that nonterminal routing rows in the
companion decision tree carry one of the four named error channels.  Once all
four normalized errors vanish on the selected terminal windows, only the five
terminal classes of \cref{table:gradient-classes-correspondence} remain.  The
scale-cascade row is included only after its raw modulation carrier has been
converted to the row-specific normalized window.  The next theorem strengthens
the classification: every entropy-critical limit is assigned to the
\emph{first} applicable class in a fixed order, eliminating ambiguity in cases
where a state could in principle satisfy more than one row of the table.

\begin{theorem}[Ordered classification of entropy-critical limits]
\label{thm:five-classes-ordered}
Every entropy-critical limit \(S_\infty\) belongs to at least one of the
following five terminal model classes.  Its \emph{ordered class} is the first
applicable class in the order below, i.e. row \(j\) with rows
\((1),\dots,(j-1)\) removed:
\begin{enumerate}[label=\textup{(\arabic*)},leftmargin=2.2em]
	\item \textbf{Compact single-core class.}
	\(\mathfrak M^{\mathrm{vel}}\) is a single localized gradient
	\(\theta\nabla\Pi\) supported in a compact subset \(K_1\subset B_R\), with
	\(K_1\) the support of one nonvanishing core; \(\Err^{*}\to0\)
	uniformly on every compact set inside \(B_R\); modulation \((a,b)\) is
	either the centered self-similar pair \((1/2,0)\) (Type~I case) or
	convergent to a constant pair (Type~II case).
	\item \textbf{Multiple-core class.}
	\(\mathfrak M^{\mathrm{vel}}\) is a finite sum
	\(\sum_{j=1}^N\theta\nabla\Pi_j\) with disjoint supports
	\(K_1,\dots,K_N\), each carrying nontrivial critical \(L^3\) mass;
	\(N\ge2\); modulation per core is bounded.
	\item \textbf{Rotational or stationary-hull class.}
	\(\partial_s V\equiv0\) modulo a constant skew-symmetric rotation
	\(\Omega\), i.e.\
	\(\mathfrak M^{\mathrm{vel}}
	=\theta\nabla\Pi+\theta(\Omega V-(\Omega y)\cdot\nabla V)\)
	for some constant \(\Omega\in\mathfrak{so}(3)\); the translation
	modulation \(b\) is not identified with \(\Omega y\); \(\Err^{*}\to0\).
	\item \textbf{Scale-cascade class.}
	\(\mathfrak M^{\mathrm{vel}}=\theta\nabla\Pi\) on each scale window
	and the modulation channel is resolved by the canonical windowwise schedule
	imported in \Cref{imp:typeii-terminal-closures}.  Before the row-specific
	window normalization this is the active carrier component \(J_1\), \(J_6\),
	or \(J_7\); on the selected entropy-critical windows the normalized
	\(\Err^{\mathrm{mod}}\) contribution vanishes.
	\item \textbf{Diffuse or tail-defect class.}
	\(\mathfrak M^{\mathrm{vel}}=\theta\nabla\Pi\) modulo a defect measure
	supported in a tail (logarithmic, Young, or coherent critical) of
	\(B_R\); the velocity converges weakly inside compacts but loses mass
	to the tail.
\end{enumerate}
The ordered classes are therefore mutually exclusive by definition, and each
entropy-critical limit is assigned to exactly one ordered class.
\end{theorem}

\begin{proof}
\textbf{Exhaustiveness.}  By \cref{def:entropy-critical}, an
entropy-critical limit \(S_\infty\) is obtained on windows where the
dissipation and all four error components vanish in the companion branch
topology.  The imported zero-gradient row-entry dictionary
\Cref{prop:conditional-row-entry-dictionary} is formulated precisely for
this terminal regime: it converts the vanishing entropy-gradient data into a
zero-gradient branch state and then assigns that state to one of the five rows
of \cref{table:gradient-classes-correspondence}.  Thus
\cref{thm:zero-gradient-recovery} places \(S_\infty\) in classes (1)--(5).

\textbf{Order assignment.}  Apply the criteria in order:

\emph{Class (1)} fires if \(\mathfrak M^{\mathrm{vel}}\) is one localized
gradient supported on a single compact subset.  This is the most
restrictive class, and its hypotheses are checked first.

\emph{Class (2)} fires if \(\mathfrak M^{\mathrm{vel}}\) decomposes into
multiple disjoint localized gradients.  This class subsumes class (1) only
in the trivial sense \(N=1\), so the raw class (2) criterion is stated with
\(N\ge2\).

\emph{Class (3)} fires if the residual symmetry data admits a constant
skew rotation \(\Omega\) but
class (1)/(2) does not apply.  Stationarity modulo rotation is
tested only after the compact-core alternatives have been removed by the
ordering.

\emph{Class (4)} fires when the modulation channel is the active carrier in
the canonical windowwise schedule before the row-specific normalization.  The
selected entropy-critical windows have normalized \(\Err^{\mathrm{mod}}\to0\),
so disjointness is determined by the imported carrier labels, not by a failure
of the entropy-critical definition.

\emph{Class (5)} fires when \(\Err^{*}\to0\) inside compact subsets but
the velocity loses mass to a tail defect.  The residual trichotomy imported in
\Cref{imp:residual-typeI-closure} routes active or affine defects back to
classes (1)--(4); class (5) is the genuinely tail-supported remainder.

\textbf{Mutual exclusion.}  Define the ordered row
\(\mathcal C_j^{\mathrm{ord}}\) as the raw row \(\mathcal C_j\) minus
\(\mathcal C_1\cup\cdots\cup\mathcal C_{j-1}\).  These ordered rows are
pairwise disjoint by construction.  Since the row-entry dictionary is
exhaustive for entropy-critical limits, every such limit lies in exactly one
ordered row.
\end{proof}

\begin{corollary}[Five-class summary table]
\label{cor:five-classes-table}
The classification of \cref{thm:five-classes-ordered} matches the
companion state space as follows.

\begin{center}\small
\setlength{\tabcolsep}{4pt}
\begin{tabular}{@{}>{\raggedright\arraybackslash}p{0.21\textwidth}
								>{\raggedright\arraybackslash}p{0.49\textwidth}
								>{\raggedright\arraybackslash}p{0.20\textwidth}@{}}
\toprule
Terminal class & Manifestations in the companion state space & Closed in \\
\midrule
(1) Compact single-core & Direct Type~I bounded ancient branches; Type~II compact single-core and retained compact-test branches & \Cref{imp:setup-branch-reduction,imp:typeii-terminal-closures} \\
(2) Multiple-core & Multibubble, separated-profile, and same-point cascade branches & \Cref{imp:typeii-terminal-closures} \\
(3) Rotational/stationary-hull & Rotational relative equilibria and stationary-hull residual branches & \Cref{imp:residual-typeI-closure} \\
(4) Scale-cascade & Scale-collapse, scale-rigid, cost-divergence, and carrier-routing branches & \Cref{imp:typeii-terminal-closures} \\
(5) Diffuse/tail-defect & Critical-tail defects, diffuse terminal defects, terminal nonactive branches, and exterior concentration & \Cref{imp:residual-typeI-closure,imp:typeii-terminal-closures} \\
\bottomrule
\end{tabular}
\end{center}
\end{corollary}

\begin{proof}
This is the manifestation column of
\cref{table:gradient-classes-correspondence}, restricted to the limiting
topology after \(\Diss\to0\) and \(\Err^{*}\to0\), with the row-specific
window normalization for class (4).
\end{proof}

\section{Reduction to the Companion Exclusion Theorems}
\label{sec:reduction}

The entropy formulation becomes effective only after it is paired with
the closure results from the companion papers.  The closure hypotheses have
already been imported in theorem-numbered form in
\Cref{imp:terminal-class-exclusion}.

\begin{theorem}[Elimination of the terminal classes]
\label{thm:eliminate-terminal-classes}
Under the admissibility, normalization, retained-mass, branch-topology, and
finite-budget hypotheses of \Cref{imp:terminal-class-exclusion}, each of the
five terminal model classes in \cref{thm:five-classes-ordered} is empty:
\begin{enumerate}[label=\textup{(\roman*)},leftmargin=2.2em]
	\item the compact single-core class is excluded by
	\Cref{imp:setup-branch-reduction,imp:typeii-terminal-closures};
	\item the multiple-core class is excluded by
	\Cref{imp:typeii-terminal-closures};
	\item the rotational or stationary-hull class is excluded by
	\Cref{imp:residual-typeI-closure};
	\item the scale-cascade class is excluded by
	\Cref{imp:typeii-terminal-closures};
	\item the diffuse or tail-defect class is excluded by
	\Cref{imp:residual-typeI-closure,imp:typeii-terminal-closures}.
\end{enumerate}
\end{theorem}

\begin{proof}
This is \Cref{imp:terminal-class-exclusion} applied to the row assigned by
\Cref{thm:five-classes-ordered}.  The role of the entropy formulation is to
show that every entropy-critical limit enters exactly one of the rows to which
the imported terminal-class exclusion package applies.
\end{proof}

Combining the assembly results, we obtain the entropy version of the
three-paper exclusion program.

\begin{corollary}[Entropy reformulation of the local exclusion program]
\label{cor:entropy-reformulation}
Assume the imported packages
\Cref{imp:setup-branch-reduction,imp:typeii-analytic-package,imp:residual-typeI-closure,imp:typeii-terminal-closures},
the exact error domination \Cref{imp:exact-error-domination}, the
zero-gradient bridge
\Cref{imp:entropy-critical-zero-gradient,prop:conditional-row-entry-dictionary},
the terminal exclusion package \Cref{imp:terminal-class-exclusion}, and the
finite corrected entropy budget hypothesis of
\Cref{prop:entropy-compactification}.  Then every entropy-critical
subsequential limit of a hypothetical singular branch is assigned to one of
the five ordered terminal classes of \cref{thm:five-classes-ordered}.  By
\cref{thm:eliminate-terminal-classes}, each such terminal class is excluded
under the stated hypotheses.  Hence no hypothetical singular branch satisfying
the imported framework and finite-budget hypotheses remains.
\end{corollary}

\begin{proof}
Combine
\cref{prop:entropy-compactification},
\cref{thm:five-classes-ordered}, and
\cref{thm:eliminate-terminal-classes}.
\end{proof}

\begin{remark}
\Cref{cor:entropy-reformulation} is an entropy-level summary of the
companion papers, not a replacement for them.  The local pressure
decomposition, carrier routing, terminal profile analysis, and residual
compactification remain the substantive PDE content.  The entropy
formalism shows that these arguments fit together as a finite
compactification mechanism for one monotone scalar quantity.
\end{remark}

\section{Comparison with Hamilton--Perelman Monotonicity}
\label{sec:hamilton-perelman}

The analogy with Hamilton and Perelman is useful but should not be
overstated.  In Ricci flow, the \(\mathcal W\)-functional is global and
geometric, with the scalar curvature appearing explicitly
\cite{Hamilton1982,Perelman2002}.  In the present Navier--Stokes setting,
there is no corresponding global geometric functional available at the
level of the companion proofs.  The pressure is nonlocal, and the branch
changes are genuinely local.  For that reason, the functional \(\Wloc\) is
localized from the outset.

Nevertheless, the two theories share a common pattern.

\begin{enumerate}[label=\textup{(\arabic*)},leftmargin=2.2em]
	\item Both use a backward adjoint weight to measure concentration near a
	candidate terminal time.  In Ricci flow, this is the conjugate heat
	kernel; here, it is the local backward adjoint weight \(\rho\) with
	modulated drift \(a y+b-V\).
	\item Both produce a monotone quantity whose defect is a sum of a
	nonnegative dissipation and explicit local error terms.  In Ricci flow,
	the dissipation is the squared norm of a gradient soliton operator; here,
	the dissipation \(\Diss_{\chi,R}\) is given explicitly by
	\eqref{eq:dissipation} and the four error terms by
	\cref{def:errors}.
	\item In both settings, a critical point of the monotone quantity
	represents a terminal model geometry.  In Ricci flow, gradient shrinking
	solitons.  Here, by
	\cref{thm:zero-gradient-recovery,thm:five-classes-ordered}, one of five
	terminal model classes.
\end{enumerate}

The decisive difference is that, in the present paper, the critical
classes are not derived from an abstract classification theorem.  They are
read off from the local state-space decomposition already established in
the companion manuscripts, and the row-by-row correspondence in
\cref{table:gradient-classes-correspondence} makes the matching explicit.

\section{Concluding Remarks}

The companion papers already provide a complete local routing
architecture.  The point of the present manuscript is to show that this
architecture admits a coherent entropy interpretation in standard PDE
language, with the calculation of \(\frac{d}{ds}\Wloc\) carried out in
full and every error term tagged to a named companion theorem.  Once the
local enstrophy, mean-free pressure, localized critical-norm velocity budget,
and backward
adjoint weight are assembled into \(\Wloc\), the decision tree may be read
as a compactness statement for the entropy flow: every singular branch
either dissipates into a previously closed local row or converges to one of
five terminal model classes.

This viewpoint suggests two further directions.  First, it separates the
purely local analytic input from the compactification mechanism, which may
be useful in formalization or in alternative singularity analyses.  In
particular, the four-piece error decomposition makes it possible to
identify \emph{which} local estimate is responsible for closing
\emph{which} branch of the architecture, a feature that should make the
proof amenable to LLM-assisted audit and formal verification.  Second, it
clarifies the role of the local pressure decomposition: the pressure does
not obstruct a monotone quantity once it is measured in the local
mean-free norms naturally produced by the companion proofs.

\subsection*{Possible extensions}

The localized entropy framework is portable to other critical PDE problems
with a similar pressure/gauge nonlocality, provided one has a local
state-space stratification on which to define the analogue of
\(\Wloc\).  Candidate settings include energy-supercritical wave
equations (with the pressure replaced by a nonlocal phase), harmonic map
heat flow with bubbling at the terminal time (with the pressure replaced
by the Lagrange multiplier of the constraint), and Yang--Mills heat flow
with the gauge as the analogue of the modulation \((a,b)\).  In each case,
the same five-row classification structure is plausible: compact
single-core, multiple-core, symmetry-reduced terminal state, scale
cascade, diffuse or tail defect.  The role of the present paper is to
provide the template; the local PDE estimates would have to be redone for
each setting.

\begin{thebibliography}{99}

\bibitem{AlbrittonBarker2019}
D.~Albritton and T.~Barker,
\emph{On local Type I singularities of the Navier--Stokes equations and
Liouville theorems},
J. Math. Fluid Mech. \textbf{21} (2019), Article 43.

\bibitem{Arnold1966}
V.~I. Arnold,
\emph{Sur la g\'eom\'etrie diff\'erentielle des groupes de {L}ie de dimension
infinie et ses applications \`a l'hydrodynamique des fluides parfaits},
Ann. Inst. Fourier (Grenoble) \textbf{16} (1966), no.~1, 319--361.

\bibitem{BallesterResidual}
G.~Duran Ballester,
\emph{Local Exclusion of Residual Type I Navier--Stokes Seregin Limits},
companion manuscript, 2026.

\bibitem{BallesterSetup}
G.~Duran Ballester,
\emph{Conditional Local Type I Navier--Stokes Reduction: Concentration,
Seregin Limits, and Liouville-Class Closures},
companion manuscript, 2026.

\bibitem{BallesterTypeII}
G.~Duran Ballester,
\emph{Type II Regularity for Navier--Stokes},
companion manuscript, 2026.

\bibitem{CalderonZygmund1952}
A.~P. Calder\'on and A.~Zygmund,
\emph{On the existence of certain singular integrals},
Acta Math. \textbf{88} (1952), 85--139.

\bibitem{CKN1982}
L.~Caffarelli, R.~Kohn, and L.~Nirenberg,
\emph{Partial regularity of suitable weak solutions of the Navier--Stokes
equations},
Comm. Pure Appl. Math. \textbf{35} (1982), no.~6, 771--831.

\bibitem{ConstantinFoias1988}
P.~Constantin and C.~Foias,
\emph{Navier--Stokes Equations},
University of Chicago Press, Chicago, 1988.

\bibitem{EbinMarsden1970}
D.~G. Ebin and J.~Marsden,
\emph{Groups of diffeomorphisms and the motion of an incompressible fluid},
Ann. of Math. (2) \textbf{92} (1970), no.~1, 102--163.

\bibitem{ESS2003}
L.~Escauriaza, G.~A. Seregin, and V.~{\v S}ver\'ak,
\emph{\(L_{3,\infty}\)-solutions of the Navier--Stokes equations and backward
uniqueness},
Russian Math. Surveys \textbf{58} (2003), no.~2, 211--250.

\bibitem{GustafsonKangTsai2007}
S.~Gustafson, K.~Kang, and T.-P.~Tsai,
\emph{Interior regularity criteria for suitable weak solutions of the
Navier--Stokes equations},
Comm. Math. Phys. \textbf{273} (2007), no.~1, 161--176.

\bibitem{Hamilton1982}
R.~S. Hamilton,
\emph{Three-manifolds with positive Ricci curvature},
J. Differential Geom. \textbf{17} (1982), no.~2, 255--306.

\bibitem{Kato1984}
T.~Kato,
\emph{Strong \(L^p\)-solutions of the Navier--Stokes equation in
\(\mathbb R^m\), with applications to weak solutions},
Math. Z. \textbf{187} (1984), no.~4, 471--480.

\bibitem{KNSS2009}
G.~Koch, N.~Nadirashvili, G.~Seregin, and V.~{\v S}ver\'ak,
\emph{Liouville theorems for the Navier--Stokes equations and applications},
Acta Math. \textbf{203} (2009), no.~1, 83--105.

\bibitem{Ladyzhenskaya1969}
O.~A. Ladyzhenskaya,
\emph{The Mathematical Theory of Viscous Incompressible Flow},
Gordon and Breach, New York, 1969.

\bibitem{LemarieRieusset2002}
P.~G. Lemari\'e-Rieusset,
\emph{Recent Developments in the Navier--Stokes Problem},
Chapman \& Hall/CRC, Boca Raton, FL, 2002.

\bibitem{Leray1934}
J.~Leray,
\emph{Sur le mouvement d'un liquide visqueux emplissant l'espace},
Acta Math. \textbf{63} (1934), no.~1, 193--248.

\bibitem{Lin1998}
F.-H.~Lin,
\emph{A new proof of the Caffarelli--Kohn--Nirenberg theorem},
Comm. Pure Appl. Math. \textbf{51} (1998), no.~3, 241--257.

\bibitem{NRS1996}
J.~Ne\v cas, M.~R\r u\v zi\v cka, and V.~{\v S}ver\'ak,
\emph{On Leray's self-similar solutions of the Navier--Stokes equations},
Acta Math. \textbf{176} (1996), no.~2, 283--294.

\bibitem{Perelman2002}
G.~Perelman,
\emph{The entropy formula for the Ricci flow and its geometric applications},
arXiv:math/0211159.

\bibitem{Scheffer1976}
V.~Scheffer,
\emph{Partial regularity of solutions to the Navier--Stokes equations},
Pacific J. Math. \textbf{66} (1976), no.~2, 535--552.

\bibitem{Seregin2012}
G.~Seregin,
\emph{A certain necessary condition of potential blow up for Navier--Stokes
equations},
Comm. Math. Phys. \textbf{312} (2012), no.~3, 833--845.

\bibitem{Stein1970}
E.~M. Stein,
\emph{Singular Integrals and Differentiability Properties of Functions},
Princeton University Press, Princeton, NJ, 1970.

\bibitem{Temam1977}
R.~Temam,
\emph{Navier--Stokes Equations: Theory and Numerical Analysis},
North-Holland, Amsterdam, 1977.

\end{thebibliography}

\end{document}
